\documentclass[11pt,a4paper]{book}

\usepackage[english]{babel}
\usepackage{csquotes}
\usepackage[T1]{fontenc}
\usepackage{lmodern}

\usepackage[a4paper,margin=2.5cm]{geometry}
\usepackage{setspace}
\onehalfspacing

\usepackage{amsmath,amssymb,amsthm}

\usepackage{graphicx}
\usepackage{booktabs}
\usepackage{caption}
\usepackage{subcaption}
\usepackage{changepage}

\usepackage{enumitem}
\usepackage{xcolor}

\usepackage[
  backend=biber,
  style=apa,
  sortcites=true,
  url=true,
  doi=true
]{biblatex}
\addbibresource{refs.bib}

\usepackage{hyperref}
\hypersetup{
  colorlinks=true,
  linkcolor=blue,
  citecolor=blue,
  urlcolor=blue,
  pdftitle={On the Indefiniteness of the Continuum Hypothesis},
  pdfauthor={Rodrigo Blanco Betes}
}

\title{\textbf{On the Indefiniteness of the Continuum Hypothesis}}
\author{Rodrigo Blanco Betes}
\date{\today}

%-------------------------
% Entornos de teoremas
%-------------------------
\theoremstyle{plain}
\newtheorem{theorem}{Theorem}
\newtheorem*{theorem*}{Theorem}
\newtheorem{lemma}{Lemma}


\begin{document}

% ---------- Portada ----------
\frontmatter
\maketitle

\tableofcontents

%-------------Introducción
\section*{Introduction}


%----------------Texto principal----------
\mainmatter

\chapter{Background}

%-------------Versiones de filosofía de las matemáticas
\section{Realism, naturalism and constructivism in philosophy of mathematics}

\begin{adjustwidth}{1cm}{1cm}
  \small
  \begin{quote}
What we find in virtually all situations is that the opposing parties ultimately defend their position by invoking extra-mathematical elements in addition to mathematical facts. For example, opponents of the ‘axiom’ of constructibility argue that it runs counter to ‘the goals of set theory’ or that it violates ‘the intended meaning of the concept of set’. What these opaque phrases suggest is that the meanings of mathematical concepts reside in the relationship in which the mind stands to the objective reality of mathematics. Moreover, if the mind is assigned a constitutive role with respect to mathematics as a whole, that relationship lies beyond the explanatory reach of mathematical facts. \parencite[p. 103]{hauser2004}
  \end{quote}
\end{adjustwidth}

\begin{quote}
So much of mathematics as is wanted for use in empirical science is for me on a par with the rest of science. Transfinite ramifications are on the same footing insofar as they come of a simplificatory rounding out, but anything further is on a par with uninterpreted systems. \textcite[p. 788]{quine1984}
\end{quote} 

Sure that intrincancy vs simplicity have not played any role at all in discussions around indispensability?
%-------------Primeros antecedentes: introducción lógica a la hipótesis del continuo
\section{The Development of CH: From a Mathematical Problem to a Philosophical Question}

%------------Exploraciones del problema más avanzadas: Woodin, MM, etc.-----
\section{The logical approach to CH after Cohen's forcing}

%-----Woodin's ideas----------------
Does the formal independence of the Continuum Hypothesis from the axioms of Set Theory imply that the question has no meaning? Some, including Cohen himself, have argued that the answer to this question is “yes” [1]. But for this position to be credible either the other classical problems in set theory which have been similarly shown to be formally unsolvable would also have to be regarded as without meaning, or there must be an explanation as to why the problem of the Continuum Hypothesis is different. \textcite[p. 29]{woodin2004}

Should the axiom of Projective Determinacy be accepted as true? One argument that is frequently cited against this axiom is that there are interesting alternative theories for the projective sets, indeed there are some very difficult open questions about the possible theories for the projective sets in the context where Projective Determinacy is false. But this argument confuses two issues. Accepting Projective Determinacy as true does not deny the study of models in which it is false. For example, it is certainly the prevailing belief that the axioms for Number Theory are consistent. Nevertheless the study of models of Number Theory in which the axioms are inconsistent seems clearly an interesting program and a deep understanding of these models probably entails resolving the P = NP question of computational complexity. Do I believe the axiom of Projective Determinacy is true? This I find to be a rather frequently asked question of me. My answer is this. I believe the axiom of Projective Determinacy is as true as the axioms of Number Theory. \textcite[pp. 31-32]{woodin2004} 

For each such countable ordinal, $\alpha$, the structure
\[
\langle \mathcal{P}(\alpha), \alpha, \cdot, +, \in \rangle
\]
is reducible to the structure,
\[
\langle \mathcal{P}(\omega), \omega, \cdot, +, \in \rangle.
\]

Therefore I claim that the next natural structure is
\[
\langle \mathcal{P}(\omega_1), \omega_1, \cdot, +, \in \rangle,
\]
which is the standard structure for all sets of countable ordinals. It is well known that
this structure is equivalent to the structure of all sets of hereditary cardinality less than $\aleph_2$.

Whether or not CH holds is a first order property of this structure. Therefore any reasonably complete axiomatization for this structure should resolve CH. \textcite[p. 32]{woodin2004}

The existence of good sentences (which are $\Omega$-satisfiable) and the theorem that good sentences must imply that the Continuum Hypothesisis false would seem to narrow the possibilities for an argument that the problem of the Continuum Hypothesisis fundamentally different from the measure problem for the projective sets to essentially just the fact that neither $CH$ or ($\neg CH$) is $\Omega_{ZFC}$-valid. But accepting this as the reason that the problem of the Continuum Hypothesis is without meaning forces one to make the stronger claim that if a 2-sentence is not $\Omega_{ZFC}$-valid then the claim that it is true is not meaningful (noting that both $CH$ and $\neg CH$ are 2 sentences). However the plausibility of such a claim critically depends on the $\Omega$ Conjecture. \textcite[p. 44]{woodin2004}

In this quote from Woodin, a sentence $\phi$ is \textit{good} in case it is $\Omega$-satisfiable, i.e. if and only if there is a consistent theory $T$ such that $T \models_\Omega \phi$. This is equivalent to saying that, for any ordinal $\alpha$ and complete Boolean algebra $\mathbb{B}$, if $V^\mathbb{B}_\alpha \models T$, then $V^\mathbb{B}_\alpha \models \phi$. This means these sentences are true in a way that is absolute and stable under forcing.

The main idea of Woodin is that in case there are “enough” good sentences, then all these sentences should lead to the fact that $CH$ is false.

%--------Proyective determinacy-----------------------

It has also been pointed out that the larger large cardinal axioms are supported by strong arguments of analogy. Typically, such arguments single out a combinatorial property of the first infinite cardinal $\omega$ and posit the existence of uncountable cardinals having an analogous property. This is usually justified by appealing to the uniform generation of sets in the cumulative hierarchy on the grounds that this uniformity should rule out the existence of ‘singularities’. It is fair to say that this line of reasoning is generally regarded as carrying less credence than justifications via reflection principles. \textcite[p. 97]{hauser2004}

Nevertheless acceptance of PD as an axiom was hampered by its lack of intrinsic evidence. That difficulty was overcome when it was shown that PD is implied by large cardinal axioms. More precisely, PD follows from the existence of infinitely many Woodin cardinals. \textcite[p. 99]{hauser2004}

Note: PD can be deduced by the existence of sufficiently large cardinal axioms. If we follow \textcite{Maddy1997} (MAXIMIZE and UNIFY), that would be an argument for accepting this axiom. However, I think that would be the case only in case fruitful conections of this axiom with more down-to-earth areas of mathematics are found. We have to consider to which extend the results on the proyective sets this axiom entails (Baire property, perfect set property, Lebesgue measurability) can be thus considered. Else, this area of research would be abbandoned through the years.

This also affords an improved understanding of the claim that PD is the ‘correct’ axiom for second order number theory. Not only does PD yield a complete structure theory for the projective sets, it is also implied by and in fact equivalent to many combinatorial principles which on the surface have nothing to do with determinacy. Most importantly, however, second order number theory as given by PD is shared by all sufficiently strong extensions of ZFC. The reason for this is the curious phenomenon that propositions implying the consistency of PD actually imply PD. \textcite[p. 100]{hauser2004}



%------------Lógica vs matemáticas. Por qué las matemáticas son sintéticas, no lógicas ni puramente formales
\section{Our stand towards logicism and the meaning of mathematical statements}

First order logic is the only that we can understand as logic. Koellner

There is not recursively ennumerable first-order theory that univocally determines the structures that mathematicians commit to.  (Gödel, Lowenheim-Skolem)

Mathematics assume to priviledge some structures, in a way thay first-order logic cannot. So mathematics ammounts to more than logic.

MacLane states that mathematics is ``formal'' in a way that could undermine these structural commitments.

When Hilbert revived Euclid’s ideas at the beginning of the twentieth century he declared this indispensability as the characteristic mark of axiomhood. In [29] [=\textcite{Hilbert1918}] for example, intrinsic plausibility is not even mentioned. (...) Hilbert’s account is inadequate in several respects. For one thing, he ignores the principle that an axiom ought to have a higher degree of evidence than the theorems it proves. Another more subtle inadequacy is brought to light by projective determanacy: how could one possibly remain convinced of its consistency without believing it is true? \textcite[p. 102]{hauser2004}

The idea of the legitimacy hierarchy does essentially provide an extrinsic reason to accept axioms as true, i.e. whether they model the most suitable structures. So essentially I agree that axioms in mathematics are justified in the end by extrinsic reasoning. In regard to the last observation by Hauser, I think that a consistent statement in a determinate axiomatic system can be regarded as a false or meaningless mathematical statement, in any case not a true mathematical statement. That does not contradict how he goes on: "The higher epistemic weight of intrinsic plausibility appears to be a consequence of the fact that—in contrast with extrinsic evidence—there is at least one way in which it is directly linked with truth, namely when an axiom is ‘implied’ by the concept(s) it aims to capture". 


%---Mathematics are formal----
\begin{adjustwidth}{1cm}{1cm}
  \small
  \begin{quote}
Moreover, there are good reasons why Mathematicians do not usually present their proofs in fully formal style. It is because proofs are not only a means to certainty, but also a means to understanding. Behind each substantial formal proof there lies an idea, or perhaps several ideas. The idea, initially perhaps tenuous, explains why the result holds. The idea becomes Mathematics only when it can be formally expressed, but that expression must be so couched as to reveal the idea ; it will not do to bury the idea under the formalism. \textcite[378]{maclane1986mathematics}
  \end{quote}
\end{adjustwidth}


%------Sobre constructivismo y convenciones--------
\section{Mathematics, constructivism and conventionalism}


%-------Sección sobre pragmática en matemáticas-----------------
\section{How pragmatic linguistics may fit into mathematics}


%----------------Primer capítulo: interpretación de CS y crítica de Koellner
\chapter{Feferman's Conceptual Structuralism and its Critics}
\label{chapter1}

\section{Introduction}
\label{Intro-Chapter-1}

At the end of the 19th century, Georg Cantor extended the concept of cardinality to infinite sets. In doing so, he was able to show that the cardinality of the continuum is strictly greater than the one of the set of all natural numbers. This led him to conjecture that there was no third set strictly between those two, in terms of cardinality; a conjecture known as the Continuum Hypothesis (CH). The same process can be extended to any set $A$, which can be shown to have smaller cardinality than its power set $\mathcal{P}(A)$. One can then formulate what is known as the Generalized Continuum Hypothesis (GCH), which states that there cannot be an intermediate set strictly between these two in terms of cardinality.\footnote{This intermediate set, in case of existing, would be called an interpolant. There is a stronger way of formulating this hypothesis in terms of well-ordered sets, but it presumes the axiom of choice; see \textcite{Koellner2023}. Even using the weaker version, CH presupposes a huge set theoretical apparatus, consisting in a system strong enough to entail the existence of a definite set containing all the subsets of the continuum, or a strong way of the power set axiom in the case of GCH.}

Cantor was unable to resolve this problem, and its status remained open for decades. Years later, \textcite{Godel1940} proved that GCH is consistent with the axioms of ZFC of set theory, and \textcite{Cohen1966} proved the same for its negation. All this gave CH its highly controversial status: an apparently well-defined mathematical question but proven to be unsolvable with the tools available to mathematicians in everyday practice.

Since then, there have been a huge number of attempts to settle CH, but none of them has been successful.\footnote{For a summary of the current status of CH, see \textcite{Feferman2011a}, Section 1. } These attempts have never offered a clear counterexample to Cantor's hypothesis, nor an intuitive model demonstrating whether it is true or false. It seems that there is no way to relate CH to a well-established and accepted result of mathematics.

This has led some mathematicians and philosophers of mathematics to doubt whether it even makes sense to seek a unique and meaningful answer to the continuum hypothesis. One of the most important proponents of this perspective is Solomon Feferman, who proposed a philosophical framework for understanding mathematics called conceptual structuralism. In this framework, mathematics is the result of open-ended processes derived from simple real-life situations, such as counting and ordering, where clarity and practical applicability are crucial. \textcite{Feferman2011a} maintained that if we accept conceptual structuralism, then CH turns into an intrinsically vague and indefinite mathematical question.

However, conceptual structuralism and its implications for CH have been subject to substantial criticism. Some of the strongest and most direct criticisms have come from Peter Koellner (\citeyear{Koellner2016,Koellner2022}). Koellner's viewpoint on CH bears similarities to those of \textcite{Parsons2016b}, Woodin (\citeyear{Woodin2001a,Woodin2001b}) or \textcite{Maddy1997}, but he distinguishes himself by having addressed Feferman's philosophical proposal more explicitly. Shortly put, Koellner's main claims are:

\begin{enumerate}
\item That conceptual structuralism is ambiguous on the status of CH.
\item That Feferman's case relies on a notion of mathematical clarity with several weaknesses, such as the fact that the metamathematical characteristics of set theory that could justify Feferman's position are also present in other areas, such as arithmetic.
\item That it is not acceptable to make mathematics dependent on subjective processes and social conventions as Feferman proposes.
\end{enumerate}

The main aim of this paper is not merely to defend Feferman against Koellner's criticisms, but to argue that the dispute itself has been framed too narrowly. Both Feferman's defenders and critics have tended to approach the determinacy of CH primarily through metamathematical and logical considerations, such as categoricity, formal semantics, and independence results. I argue that these considerations are insufficient on their own. Questions concerning mathematical determinacy must also take into account the role different mathematical structures play within scientific and mathematical practice. From this perspective, arithmetic and higher set theory occupy importantly different positions, and these differences help explain why CH appears philosophically less determinate than ordinary arithmetic truths.

A proper treatment of point (3) above would require a deep examination of speech act theory and institutional theory; something that is well beyond the scope of this paper. Let me briefly point out, however, that I do not think, as Koellner does, that it is impossible to understand mathematical entities in an institutional or social way. Quite the contrary. There have been, for the time being, promising attempts to understand mathematics in an institutional way, as Feferman does, and they have proven to be quite resilient to criticism.\footnote{Take for example \textcite{Cole2013}, who gives an institutional account of mathematics based on speech act theory and argues that it can give an account of the necessity and objectivity of mathematical facts.}

In summary, the most important thesis this paper is that an alternative interpretation can be given to conceptual structuralism, which focuses most on how mathematical structures are approached and applied in real-life practice, and less on logical and epistemological clarity. I then argue that the strongest version of Feferman's position does not rest primarily on metamathematical considerations, but on the role different structures play within mathematical and scientific practice.

The text is organized as follows: I first reconstruct Feferman's conceptual structuralism and its relation to CH, in Section \ref{conceptual-structuralism-intro}. I then examine the role of scientific applicability and mathematical practice in evaluating the determinacy of mathematical structures, arguing that these considerations complement Feferman's notion of conceptual clarity. In Sections \ref{categoricity-second-order-logic} and \ref{standard-models}, I discuss categoricity, second-order logic, and the problem of standard models, comparing the cases of arithmetic and set theory. 

Finally, I argue that the asymmetries between these domains in ordinary mathematical practice support a revised interpretation of conceptual structuralism under which the Continuum Hypothesis may reasonably be regarded as indefinite.



\section{Understanding Feferman's conceptual structuralism}
\label{conceptual-structuralism-intro}

It is not always easy to follow the essential claims of Feferman's conceptual structuralism.\footnote{Other approaches to the philosophy of mathematics have also been characterized as \textit{conceptual structuralism}, such as that of \textcite[postnote]{Isaacson2008}. Unless otherwise stated, however, references to conceptual structuralism in this text concern Feferman's version of the view.} The fundamental idea of this conception is that mathematical notions come from structures based on processes of real experience of activities such as counting, ordering, joining, combining, etc.; and that it is within these structures that mathematical objects and systems are derived and constructed. However, \textcite{Feferman2011a} characterizes it in ten theses that are arguably overly concise. He himself acknowledged that they would require further elaboration, but he did not have the space to provide it:

\begin{adjustwidth}{1cm}{1cm}
\begin{quote}
\small
There is no time to elaborate these points, but various aspects of them will come up in our discussion of two constellations of structural notions, first of objects generated by one or more “successor” operations, and second of the continuum. \textcite[p. 12]{Feferman2011a}.
\end{quote}
\end{adjustwidth}

This concerns us here because the most direct criticism by Koellner of Feferman's argument is precisely that his philosophical conception does not say anything about whether CH must be considered a definite statement:

\begin{adjustwidth}{1cm}{1cm}
\begin{quote}
\small
The main point I wish to stress is that conceptual structuralism does not on its own say anything about whether or not a given conception (the natural numbers, arbitrary sets of natural numbers) leads to definite or indefinite statements. We are thus provided with a flexible framework, one that is capable of being applied to views on which CH is definite and also to views on which CH is indefinite. \textcite[p. 10]{Koellner2022}.
\end{quote}
\end{adjustwidth}

%------------------------------
I think that Feferman's writings on the topic were more successful in making his position explicit than in defending it. As we will see in this section, he argues, both philosophically and formally, that mathematical conceptions must be clear and conceptually and epistemologically grounded. However, he is less explicit about why this must be so, and about why conceptions such as the one of the set-theoretical hierarchy fail to satisfy these requirements. I think this gap could be addressed through a closer evaluation of mathematical practice, a point to which I return in Section \ref{science-applicability}.

A central idea in conceptual structuralism is that mathematical conceptions always refer to mathematical structures and must be grounded in simple and elementary processes such as counting, combining, and joining. A recurring theme in Feferman's work is that these processes should correspond closely to human sensorimotor and cognitive capacities, a point emphasized in Thesis 3 of conceptual structuralism:

\begin{adjustwidth}{1cm}{1cm}
\begin{quote}
\small
The basic conceptions of mathematics are of certain kinds of relatively simple ideal world-pictures which are not of objects in isolation but of structures, i.e. coherently conceived groups of objects interconnected by a few simple relations and operations. They are communicated and understood prior to any axiomatics, indeed prior to any systematic logical development. \textcite[p. 11]{Feferman2011a}.
\end{quote}
\end{adjustwidth}

Feferman is likely thinking here of basic conceptions such as the notion of natural a number, the induction principle, and a pre-theoretical idea of the successor operation. From these, some further logical development may be appropriate.\footnote{In this regard, the idea of unfolding is relevant. Unfolding refers to taking a concept, definition, or system and systematically developing its consequences, rules, or elements in a transparent, step-by-step manner. See \textcite{FefermanStrahm2000}.} These basic conceptions then become part of an ideal world: a coherent conceptual scenario adopted as a framework for doing mathematics. They are “ideal” because they abstract away from physical reality, yet they remain “world pictures” insofar as they closely correspond to basic real-life experiences. \footnote{It is interesting to note that, in formulating this thesis, Feferman refers to “ideal world-pictures”, linking the last two words, whereas when discussing the conception of the set-theoretical hierarchy—see, for instance, \textcite[p.~19]{Feferman2011a}—he speaks of an “ideal-world picture”. In \textcite[p.~19]{Feferman2011b}, however, he writes “ideal world pictures”, without a hyphen, which arguably aligns more closely with the first of these interpretations.} Feferman also states here that a mathematical conception must be a conception of a structure, a view that I adopt throughout the rest of this text. A distinction is sometimes drawn between “monomorphic” structures, which are categorically determined, and “polymorphic” ones, which are not—see \textcite[133]{Ferreiros2023ConStructuralism}. This distinction is closely related to Feferman's notion of clarity, though the two are not necessarily interchangeable, as I explain in \ref{categoricity-second-order-logic}.

Through this notion, Feferman distinguishes between clear and less clear conceptions of mathematical entities: a clear conception allows statements about it to have determinate truth values, whereas a less clear one can give rise to vague or indeterminate statements, such as the Continuum Hypothesis. This argument integrates naturally with the maxims of conceptual structuralism, since Thesis 5 states:

\begin{adjustwidth}{1cm}{1cm}
\begin{quote}
\small
Basic conceptions differ in their degree of clarity. One may speak of what is true in a given conception, but that notion of truth may only be partial. Truth in full is applicable only to completely clear conceptions. \textcite[p. 11]{Feferman2011a}.
\end{quote}
\end{adjustwidth}

Here, Feferman explicitly acknowledges that truth is not absolute for all conceptions---only fully clear ones can support determinate truth claims. CH, as a statement arising from less clear set-theoretic conceptions, exemplifies this indeterminacy.

%-----------------------------------------
This philosophical approach, presented in Section 2 of \parencite{Feferman2011a}, has a formal counterpart in Section 3. There, Feferman introduces a method for determining whether a mathematical statement is well defined, building on his earlier work on semi-constructive theories \parencite{Feferman2010}. He argues that the Continuum Hypothesis may ultimately prove to be an indefinite statement. In these papers, Feferman adopts an intermediate position between classical and intuitionistic logic, using classical logic in contexts he considers appropriate while relying on intuitionistic logic elsewhere.
Again, I take this logical approach to CH not to provide a new argument, but rather to offer an explanation. Feferman essentially reformulates his philosophical position in formal terms, sometimes making more explicit some of its otherwise obscure aspects. However, I do not think that this formal framework provides any additional reason to accept his position for those who are unconvinced by the underlying philosophical arguments.

Feferman develops the formal system $\mathrm{SCS}$ (Semi-Constructive Set Theory), based on an intuitionistic version of the Kripke–Platek (KP) axiomatic system together with several additional principles. Schematically, the system can be represented as:

\begin{equation}
\mathrm{SCS} = \mathrm{IKP} + \mathrm{MP} + \mathrm{AC}
\end{equation}

Here, $\mathrm{IKP}$ denotes the intuitionistic Kripke–Platek system with the Axiom of Infinity, $\mathrm{MP}$ denotes Markov's Principle, and $\mathrm{AC}$ denotes a constructive and weaker version of the Axiom of Choice.\footnote{See \textcite[Section 6]{Feferman2010} for further details regarding the added axioms. The system defined there also includes the bounded omniscience scheme and the law of excluded middle for $\Delta_0$ formulae, but \textcite{Rathjen2016} proved these principles to be redundant.}

According to Feferman, the validity of the Law of Excluded Middle, $A \lor \neg A$, for a given formula $A$ indicates that $A$ is definite; similarly, other classical principles are associated with definiteness. $\mathrm{SCS}$ assumes definiteness for $\Delta_0$ formulas and for first-order arithmetical statements. By contrast, the universe of sets as a whole is regarded as indefinite, which explains why the system generally operates within intuitionistic logic.

Feferman subsequently conjectured that CH is indefinite in this axiomatic system; that is, that one cannot derive $\mathrm{CH} \lor \neg \mathrm{CH}$. He went further still, conjecturing that CH also remains indefinite in the stronger system $\mathrm{SCS} + \mathcal{P}(\mathbb{N})$, which postulates the existence of the power set of the natural numbers. Both conjectures were later confirmed by \textcite{Rathjen2016}.

Rathjen, however, was cautious about drawing philosophical conclusions from the result and emphasized instead its technical character:

\begin{adjustwidth}{1cm}{1cm}
\begin{quote}
\small
The objective of this paper is to prove Feferman's conjecture. In this sense it is a technical paper (...). However, as far as the ongoing discussions of the foundational status of CH are concerned, readers will have to form their own conclusions. \textcite[p. 743]{Rathjen2016}.
\end{quote}
\end{adjustwidth}

A similar attitude is expressed by \textcite[p. 516]{Koellner2016}: “It seems antecedently unsurprising that if one assumes only classical logic for a limited domain—say first-order number theory—then one will not be able to prove that classical logic holds for distinctive statements of a richer domain.” This fits well with the idea that the formal approach serves primarily to articulate more precisely what counts as “clear” within conceptual structuralism, rather than to provide an independent justification for Feferman's position. However, unlike Koellner, I am inclined to think that Feferman himself never intended this formal approach to serve a stronger justificatory role, or at least never claimed that explicitly.

This formal approach also reveals a tendency toward predicativism, particularly in its rejection of the full power set operation on the set of natural numbers. In SCS, real numbers are constructed via Cauchy sequences rather than as arbitrary subsets of $\mathbb{N}$. While this suffices for much of analysis, it avoids commitment to the full set-theoretic conception of the continuum.

Feferman nevertheless acknowledges a more moderate position, according to which $\mathcal{P}(\mathbb{N})$ may be accepted while higher levels of the set-theoretic hierarchy are rejected. Importantly, his formal results concerning CH remain valid even under this more permissive stance, since CH remains indefinite in the stronger system $\mathrm{SCS} + \mathcal{P}(\mathbb{N})$. This also fits well with Feferman's broader philosophical outlook, as can be seen, for instance, in (\citeyear[p.~22]{Feferman2014ConceptualStruc}):

\begin{adjustwidth}{1cm}{1cm}
\begin{quote}
\small
The direct apprehension of these [structures] leads one to speak of truth in a structure in a way that may be accepted uncritically when the structure is such as the integers but \textit{may} be put into question when the conception of the structure is less definite as in the case of the geometrical plane or the continuum, and \textit{should} be put into question when it comes to the universe of sets.
\end{quote}
\end{adjustwidth}


\section{Scientific applicability in evaluating Feferman's case}
\label{science-applicability}

We have seen that Feferman imposes important epistemological limitations on mathematical structures, at least if statements formulated within them are to be regarded as having determinate truth values. These limitations are relatively clear—especially when expressed formally through the systems $\mathrm{SCS}$ and $\mathrm{SCS} + Pow(\omega)$—but it is far less clear how they are to be evaluated and ultimately justified on philosophical grounds.

As I explain later in Sections \ref{categoricity-second-order-logic} and \ref{standard-models}, such claims are difficult to justify purely in formal or metamathematical terms. For this reason, I would like to explore a somewhat different approach. Rather than focusing exclusively on formal considerations, this approach asks to what extent the axiomatic systems and structures that would render CH determinate are integrated into scientific and mathematical practice.

This does not mean, of course, that we should naïvely assume that every mathematical structure or theory employed in practice is thereby determinate, nor that all statements formulated within such structures possess determinate truth values. Nevertheless, the integration of a mathematical structure into mathematical practice usually involves at least some degree of commitment to it on the part of practitioners. Full commitment to a structure—that is, a platonistic belief in its reality—would, of course, imply that statements formulated within it are determinate.

For this reason, the nature of scientific and mathematical practice, and the way in which different mathematical structures become integrated into those practices, may play an important role in debates concerning the determinacy of CH. Apart from this, it is also important to understand the extent to which Feferman himself takes mathematical structures to be integrated into practice, beginning with scientific practice outside mathematics. A corresponding discussion of mathematical practice itself is deferred to Section \ref{mathematical-practice}.

Feferman was generally skeptical of the idea that mathematical entities exist in an abstract Platonic realm independent of reality. At the same time, he acknowledged that the laws of physics are formulated in highly idealized terms, and that scientific models do not necessarily correspond directly to the real world:

\begin{adjustwidth}{1cm}{1cm}
\begin{quote}
\small
One must ask whether the mathematical structure of the real number system can be identified with the physical structure, or whether it is instead simply an idealized mathematical model of the latter, much as the laws of physics formulated in mathematical terms are highly idealized models of aspects of physical reality. \textcite[p. 107]{Feferman1999}.
\end{quote}
\end{adjustwidth}

Koellner even regards this as a difficulty for Feferman's position, observing:

\begin{adjustwidth}{1cm}{1cm}
\begin{quote}
\small
There is a hint of an answer to this question in his remarks on the role of mathematical models in physics. For in addition to thinking that no mythical third realm (the platonic realm) can secure the definiteness of set theory, he also thinks that the physical world cannot secure the definiteness of set theory. \textcite[p. 507]{Koellner2016}.
\end{quote}
\end{adjustwidth}

However, the fact that the physical world does not correspond directly to any mathematical structure does not imply that idealized mathematical conceptions are entirely independent of physical reality. What seems most consistent with Feferman's work is instead the view that mathematical entities are justified pragmatically rather than ontologically. On this interpretation, the crucial link between mathematics and reality lies in our practices—in the ways we model, measure, and reason about the world.

This line of thought bears some resemblance to the indispensability arguments of Quine and Putnam, which are noteworthy insofar as they were among the first to emphasize the role of the applicability of mathematics to the world in justifying belief in mathematical entities. Feferman's position, however, differs from theirs in important respects. Whereas Quine and Putnam were primarily concerned with epistemological and ontological questions, Feferman places much greater emphasis on mathematical practice and on the epistemological limitations governing mathematical conceptions.\footnote{Indeed indispensability argument are far from universally accepted among philosophers of mathematics. For a critical discussion, see \textcite{Maddy1992}.}

For this reason, even without assuming a complete correspondence between mathematical entities and physical objects, there remains room to maintain that, within conceptual structuralism, mathematical structures may bear an indirect relation to physical reality through processes of modelling, idealization, and abstraction.

Feferman conjectured that all mathematics applicable to real-world physics is derivable from relatively simple systems such as Peano Arithmetic. If so, systems such as $\mathrm{SCS}$—which reproduce the formal assumptions of conceptual structuralism as an account of mathematics—would suffice to formalize all scientifically relevant mathematical applications without requiring stronger set-theoretic frameworks:

\begin{adjustwidth}{1cm}{1cm}
\begin{quote}
\small
I have come to conjecture that practically all scientifically applicable mathematics can be formalized in systems reducible to PA, or, as I have sloganized it in \textcite{Feferman1993}: a little bit goes a long way. Against this, I have learned of a couple of cases in some approaches to the foundations of quantum field theory where it appears one must go beyond the resources of PA; but the physical theories that require such additional strength are rather speculative. \textcite[p. 15]{Feferman1999}.
\end{quote}
\end{adjustwidth}

This remains, of course, a conjecture, but it is a fairly plausible one. Most physical theories are formulated within the frameworks of special relativity, quantum mechanics, or general relativity, and none of these appears to require especially strong foundational systems. Even if the claim that all such mathematics can be formalized in PA is too strong, stronger but still comparatively weak systems such as $\mathsf{ACA}_0$ remain available.

The system $\mathsf{ACA}_0$ can be formulated within second-order arithmetic and permits the formation of sets of natural numbers definable by arithmetical formulas. Calculus, algebra, functional analysis, basic topology, and many other branches of mathematics used in physics can be formalized within this framework. The stronger system $\Pi^1_1\text{-}\mathsf{CA}_0$, which additionally allows sets of natural numbers defined by $\Pi^1_1$ formulas, is (thought not predicative) generally regarded as more than sufficient. In any case, there is no need to move beyond the language of second-order arithmetic.\footnote{These systems arise in the program of reverse mathematics, which is largely carried out within second-order arithmetic. See \textcite{Simpson2006} for further details.}

By contrast, statements such as the Continuum Hypothesis require systems of third-order arithmetic, quantifying not only over natural numbers and sets of natural numbers (that is, real numbers), but also over sets of sets of natural numbers. This highlights the fact that merely formulating CH already involves a level of set-theoretic complexity far beyond what is required for the core mathematics used in physics. Thus, at least from the standpoint of scientific applicability, there appears to be little reason to regard CH as determinate, given that scientific practice would remain essentially unaffected regardless of whether CH were true or false. Needless to say, scientific applicability is not the only consideration that might support regarding CH as determinate.

This is another point at which Feferman's tendency toward predicativism appears, since he argues \parencite[p.~22]{Feferman2011a} that all current applications of analysis in physics can be carried out predicatively, whithout need of set-theoretically defining the continuum:

\begin{adjustwidth}{1cm}{1cm}
\begin{quote}
\small
What this shows is that the argument from the success of physics to the reality of the set-theoretical conception of the real numbers is completely vitiated.
\end{quote}
\end{adjustwidth}

One might nevertheless argue that this reduces the scientific applicability of mathematical conceptions merely to their capacity to provide the formal tools required for proving relevant results in physics. Scientific and mathematical practice involve more than the development of frameworks within which results can be formulated and proved; they constitute networks of interrelated practices involving, for example, practical applications, communication, pedagogy, and the dissemination of mathematical knowledge.

From this perspective, predicative mathematics may sometimes complicate the construction and development of certain mathematical domains, especially regarding the problem of constructing the continuum in a predicative and non-set-theoretic manner. Still, I think Feferman was right to remain flexible on this issue, as suggested at the end of Section \ref{conceptual-structuralism-intro}. Moreover, his position concerning CH does not appear to be significantly affected here, since the limits of predicative mathematics are generally located within the scope of second-order arithmetic. To assess this issue more fully, however, we must move beyond scientific applicability and consider mathematical practice itself.

\section{Conceptual Structuralism regarding the philosophy of mathematical practice}
\label{mathematical-practice}

For many years, the philosophy of mathematics was dominated by the foundationalist tradition, closely associated with foundational programs within mathematics, such as axiomatic set theory. This situation began to change in the second half of the twentieth century with the emergence of the so-called \textit{philosophy of mathematical practice}. As \textcite{Mancosu2008} observes:

\begin{adjustwidth}{1cm}{1cm}
\begin{quote}
\small
Mathematical logic, which had been essential in the development of the foundationalist programs, was seen as ineffective in dealing with the questions concerning the dynamics of mathematical discovery and the historical development of mathematics itself. \textcite[p. 4]{Mancosu2008}.
\end{quote}
\end{adjustwidth}

Among the central ideas commonly associated with this philosophical tradition are the following:\footnote{In our discussion of mathematical practice, we follow in several respects the way \textcite{Ferreiros2016} approaches the concept; see especially Section 2.3.}

\begin{enumerate}
\item Mathematics is a human activity, not merely a compendium of timeless truths. In this activity, reference to the mathematical community is essential.
\item Explanations, clarity, and heuristics play an important role in mathematics, and not only formal reasoning.
\item Social and historical elements play a crucial role in the development of mathematics. There is not necessarily a unique framework that univocally determines mathematical practice.
\item Theoretical knowledge is essential to mathematics, but must ultimately remain connected to practice.
\end{enumerate}

Regardless of whether one accepts the broader commitments of this philosophical movement, it is easier to agree that any principle consistently guiding mathematical practice constitutes an important datum that the philosophy of mathematics must account for. For this reason, and in addition to the considerations discussed in Section \ref{science-applicability}, it is important to understand how Conceptual Structuralism relates to this perspective, and what role structures in which CH is regarded as determinate play within mathematical practice more generally.

%----------------Maddy's MAXIMICE&UNIFY-----------------------
One of the most influential defenses of methodological principles guiding mathematical practice is that developed by \textcite{Maddy1997}, who emphasizes two complementary maxims: \emph{MAXIMIZE} and \emph{UNIFY}. According to this view, mathematical practice tends to expand the mathematical universe as much as possible, provided consistency is preserved. This resembles the principle of quasi-combinatorialism, that is, the idea that understanding set theory involves grasping the notion of an arbitrary set—definable or not—and therefore accepting as many sets as possible.

Such principles were undoubtedly important in the early development of set theory, but they have not remained uniformly dominant, either prior to the rise of foundationalism or in contemporary mathematical practice. Moreover, maximizing principles may conflict with one another; consider, for example, the incompatibility between the Axiom of Choice and the Axiom of Determinacy. Furthermore, the very notion of an arbitrary set is itself controversial, partly because it cannot be fully captured within formal systems such as ZFC, or even within considerably stronger systems, as argued by \textcite{Ferreiros2011}.

%----------Feferman & mathematical practice-----------------
Feferman's conceptual structuralism can in fact be understood as an alternative to unrestricted maximizing and unifying principles; see \parencite[p.~9]{Feferman2011a}. He addresses Maddy's proposal more directly in \textcite[p.~9]{Feferman2000NewAxioms}:

\begin{adjustwidth}{1cm}{1cm}
\small
\begin{quote}
In binding itself to mathematical practice, this kind of naturalism is in danger of being unduly transitory. Even if one takes the proposed naturalistic point of view and mathematical practice as exemplified in set theory for granted, there is a crucial question as to what determines the mathematical ends'' for which the most effective means'' are to be sought. And having chosen the ends, in what sense does effectiveness justify the means?
\end{quote}
\end{adjustwidth}

On the same page, Feferman distances himself further from Maddy’s maxims by arguing that the mathematical universe should, in some sense, remain limited, even if no absolute or rigid limiting principle can be imposed:

\begin{adjustwidth}{1cm}{1cm}
\small
\begin{quote}
Old-style infinitesimals, Dirac delta-functions, unrestricted comprehension? If not, what justifies what is to be admitted to mathematics? Once it is agreed that there has to be some sort of justification, intrinsic or extrinsic, then one is in the game of potentially limiting mathematics in some way or other. I don't hew to any sort of absolute principle in favor of limiting mathematics.
\end{quote}
\end{adjustwidth}

In these passages, Feferman suggests that merely following whatever mathematical practitioners happen to do at a given historical moment may lead to philosophically fragile accounts of mathematics. At the same time, he is not claiming that the philosophy of mathematics should ignore mathematical practice altogether. Rather, he implicitly appeals to mathematical practice in arguing against maximizing conceptions of mathematics, noting that entities or principles such as infinitesimals, Dirac delta-functions, and unrestricted comprehension are no longer accepted in the unrestricted forms in which they originally appeared. In my view, the key is instead to identify methodological principles that remain stable across historical periods and across different domains of mathematical practice.

The question, then, is not whether mathematical practice matters, but which features of mathematical practice could justify regarding certain mathematical structures as determinate. In some cases, the role a structure plays within mathematical practice gives us stronger reasons to regard it as determinate than can be provided by metamathematical results alone. In particular, a structure such as the set $\mathbb{N}$ of natural numbers is treated significantly differently in mathematical practice from structures such as the set-theoretic hierarchy. The following features do not by themselves settle questions of determinacy, but they help explain why certain mathematical structures are treated in practice as more fixed and objective than others:

\begin{enumerate}

  \item \textbf{Modes of access in practice:} Our practical interaction with the natural numbers is iterative, sequential, and intuitive. Some of our most fundamental cognitive capacities, especially those related to space and time, appear to mirror this structure. By contrast, set-theoretic conceptions appear quite differently in mathematical practice: they are more dependent on formal axiom systems and less directly connected to intuition. For instance, ordinary mathematical practice almost never directly manipulates objects such as $\mathcal{P}(\mathbb{R})$.\footnote{If we examine the paradigmatic structures that arise across different areas of mathematics, we find that they often remain below the level of the full power set of the continuum. For example, the canonical structure in topology is the real line endowed with its usual topology, which restricts attention to open sets. Similar observations apply in areas such as geometry, analysis, and measure theory.}

  \item \textbf{Degree of framework dependence:} Arithmetic often functions as a fixed background across different foundational programs. Apart from finitists and certain forms of constructivism and intuitionism, most foundational perspectives accept the basic notions of arithmetic and treat arithmetical statements as determinate. By contrast, higher set theory is often highly sensitive to variations in the surrounding foundational assumptions. This is reflected in the large number of competing axiomatic frameworks that have been proposed for set theory, including $\mathrm{ZFC}$, $KP$, intuitionistic set theory, $\mathrm{SCS}$, and many other cases.
  
  \item \textbf{Explanatory indispensability vs. optional enrichment:} The natural numbers seem unavoidable in ordinary mathematics. By contrast, large portions of higher set theory function more like optional enrichments of mathematics. The Continuum Hypothesis is a paradigmatic example: vast areas of mathematics proceed unaffected regardless of its truth value. Similar considerations apply to large cardinal axioms, determinacy principles, forcing axioms, and related set-theoretic extensions.
  
  \item \textbf{Convergence of representations:} The natural numbers appear more determinate in part because an enormous variety of mathematical practices converge on the same structure. By contrast, the set-theoretic hierarchy admits radically different extensions and interpretations, including forcing extensions, inner models, large cardinal hierarchies, determinacy models, and multiverse conceptions.
  
\end{enumerate}

In Section \ref{standard-models}, I return to these differences, focusing in particular on how mathematical practice treats the standard models of arithmetic and set theory. In the context of Conceptual Structuralism, it is natural to view these features of mathematical practice as reflecting one of its central ideas: namely, that determinate structures must remain sufficiently connected to our cognitive capacities. I am not claiming that this automatically proves Feferman right, but rather that a sufficiently detailed examination of mathematical practice may have something to contribute to the question of the definiteness of CH.

I now return to metamathematical arguments, which play an important role in Koellner's criticism. I agree with Koellner that such arguments are not sufficient to establish the indefiniteness of CH, but I also think that he assigns them too much weight in his interpretation of Feferman.


\section{Categoricity and second order logic}
\label{categoricity-second-order-logic}

One of the main metamathematical and model-theoretical approaches to determinacy of mathematical structures is the concept of categoricity. Perhaps the strongest reason for regarding a mathematical structure as determinate would be the existence of a categorical axiomatization in first-order logic. However, the Löwenheim--Skolem theorem implies that no recursively enumerable first-order theory with an infinite model can be categorical. Consequently, neither the structure of the natural numbers nor the set-theoretic hierarchy can be regarded as determinate solely on the basis of first-order categoricity.

Some authors, such as \textcite[31]{Isaacson2008} and \textcite[chapter 8]{shapiro1991foundations}, have argued that categoricity results in full second-order logic should be accepted, since categoricity becomes available in that setting. In the case of arithmetic, full second-order Peano arithmetic, $PA_2$, is categorical. In the case of set theory, an initial segment of the universe $V_\kappa$, where $\kappa$ is a strongly inaccessible cardinal, is likewise quasi-categorically characterized by second-order formulations of Zermelo-Fraenkel set theory, $\mathrm{ZFC}_2$. Both \textcite[19]{Feferman2011a} and \textcite[511]{Koellner2016} reject this position, arguing that full second-order semantics presupposes a background notion of ``all subsets,'' thereby involving a commitment of essentially the same kind and strength as the set-theoretic ontology under discussion. 

In this section, I sketch both the categoricity result for arithmetic and the quasi-categoricity result for set theory. I then examine whether, even if these results fail to establish the determinacy of the corresponding structures, they nevertheless receive different treatment within ordinary mathematical practice.

\subsection{The categoricity of \texorpdfstring{$\mathrm{PA}_2$}{PA2}}

The categoricity theorem for second-order Peano Arithmetic ($\mathrm{PA}_2$) can be formalized within $\mathrm{ZF\!-\!P}+\mathrm{Pow}(\omega)$. Under full second-order semantics, second-order variables range over all subsets of the first-order domain. In the case of arithmetic, this amounts to quantification over $\mathcal{P}(\mathbb{N})$. The crucial principle is the second-order induction axiom:

\[
\forall X\Bigl[(0\in X \wedge \forall n\,(n\in X \rightarrow S(n)\in X))
\rightarrow \forall n\,(n\in X)\Bigr].
\]

Once this principle is assumed together with the remaining axioms of $PA_2$, the proof that every model of the theory is isomorphic to the standard model is straightforward.  Be $(\mathbb{N},0,S)$ the standard model of $PA_2$, and suppose that $(M,0_M, S_M)$ is another structure that models this system of axioms. Let $(\mathbb{N},0,S)$ denote the standard model of $PA_2$, and suppose that

\[
(M,0_M,S_M)\models PA_2.
\]

Define a function $f:\mathbb{N}\to M$ recursively by

\[
f(0)=0_M,
\qquad
f(n+1)=S_M(f(n)).
\]

The successor axioms imply that $f$ is injective. To prove surjectivity, define $A=\{x\in M:\exists n\in\mathbb{N}\,(f(n)=x)\}$. By construction, $0_M\in A$. Moreover, if $x\in A$, then for some $n\in\mathbb{N}$ we have $x=f(n)$, and hence $S_M(x) = S_M(f(n)) = f(n+1)$, so $S_M(x)\in A$. Thus $A$ contains $0_M$ and is closed under $S_M$. By the second-order induction axiom, $
A=M$.

Hence every element of $M$ lies in the image of $f$, so $f$ is surjective. Therefore $f$ is a bijection preserving $0$ and successor, and thus an isomorphism. Consequently every full model of $PA_2$ is isomorphic to $(\mathbb{N},0,S)$.

I have highlighted only those aspects of the proof that are relevant to the assumptions on which it depends, rather than presenting its full formal derivation.\footnote{For a rigorous formulation of the proof, see \textcite[82]{shapiro1991foundations}.}

The crucial point, in my view, is that the underlying set-theoretic assumptions are comparatively weak. The strongest assumptions involved are arguably the existence of the power set of the natural numbers together with a weak form of replacement, both of which tend to be accepted ubiquitously throughout ordinary mathematical practice.

\subsection{The quasi-categoricity of ZFC2}

The quasi-categoricity result of $\mathrm{ZFC}_2$ can be formalized and proven within $\mathrm{ZFC}$, a similar setting that we found in arithmetic. We can formalize in this system the idea of being a model of full second-order set theory by the relation $(M,P) \models \mathrm{ZFC}_2$, where $P=\mathcal{P}(M)$ (and similarly for the interpretation of higher order variables), and in such a way that the $\mathrm{ZFC}$ axioms hold when relativized to $M$, where the replacement schema can now be seen as a single axiom:

\[
\forall f\,\forall x\,\exists y\,\forall z\,
\Bigl(
z\in y \leftrightarrow \exists t\in x\,(z=f(t))
\Bigr).
\]

With this setting, the quasi-categoricity result is stated as follows:

\begin{theorem*}[\textcite{Zermelo1930}]
  
\(\kappa\) is strongly inaccessible iff
\[
V_\kappa \models \mathrm{ZFC}_2 .
\]

\end{theorem*}

\begin{proof}

Except replacement, the axioms of $\mathrm{ZFC}$ easily hold in \(V_\alpha\), in case \(\alpha > \omega\) is a limit ordinal. Replacement is the axiom that uses the hypothesis that $\kappa$ is inaccessible: if $x \in \kappa$ and $f: x \to \kappa$, $f(x)= \{ y \subseteq \kappa | \exists z \in x, f(z) = y \}$, then we have $|f(x)| \leq |x| < \kappa$. As this cardinal is inaccessible, it cannot be reached by taking unions or power sets of smaller cardinals, which implies that $f(x) \in V_\kappa$, thus being replacement verifyed relatively to this model.

To proof the converse, we need to see both that $\kappa$ is regular and that, if $\mu < \kappa$, $\mathcal{P}(\mu) < \kappa$. If $\kappa$ were not regular there would be an $\alpha < \kappa$ and a function $f : \alpha \to \kappa$ such that $max_{x\in\alpha}\{ f(x) \} = \kappa$. Then $f \subseteq V_\kappa$, so by second-order Replacement $Im(f) \in V_\kappa$. Then by the axiom of union in $V_\kappa$, $\cup Im(f) = \kappa \in V_\kappa$, a contradiction. For the second part, suppose there would be a \(\mu < \kappa\) such that $2^\mu \geq \kappa$. Then there would be a surjective function $f : \mathcal{P}(\mu)\to \kappa$, But power set axiom implies $\mathcal{P}(\mu) \in V_\kappa$ so by second-order replacement, $Im(f)= f(\mathcal{P}(\mu))=\kappa \in V_\kappa$, a contradiction.\footnote{See \textcite[23]{Kanamori2003} for a modern formulation and proof of the theorem.}

\end{proof}

Note that, similarly to the case of arithmetic, the proof is straightforward. The theorem then shows that $V_\kappa$ is the only model of set theory of the corresponding size, what leads to the \textit{quasi-categoricity} label. 

\subsection{Comparison}

Note that, in the end, this is a theorem of first-order $\mathrm{ZFC}$. More precisely, $\mathrm{ZFC}_1$ formalizes and proves that for every strongly inaccessible cardinal $\kappa$, $V_\kappa$ is the unique full second-order model of $\mathrm{ZFC}_2$ of cardinality $\kappa$. Internally, the relevant $V_\kappa$ are isomorphic (i.e., within a given model, $V_\kappa$ is trivially isomorphic to itself). Externally, however, two models of $\mathrm{ZFC}$ may have different $V_\kappa$'s.

The situation in arithmetic is parallel. Working in $\mathrm{ZF\!-\!P}+\mathrm{Pow}(\omega)$, one can formalize the relation $\mathcal{M} \models \mathrm{PA}_2$ and prove the corresponding categoricity theorem. However, $\mathrm{ZF\!-\!P}+\mathrm{Pow}(\omega)$ is itself not categorical, and the corresponding structures of natural numbers in two different models of the theory need not be isomorphic.

So it seems that both results exhibit a potential circularity, since they do not establish the determinacy of the intended structures unless the higher-order metatheoretical structures in which these results are formulated and proved are themselves already taken to be determinate. In this vein, \textcite[20]{Feferman2011a} states that ``the circularity is so patent, one does not see that any more can be claimed thereby than that the definiteness of those notions may be accepted as a matter of faith.''

In either case, this does not provide a means of establishing definiteness in purely logical terms, at least not within first-order finitary logic, which has classically been regarded as the canonical notion of logic. If one instead adopts full second-order logic, then the relevant determinacy assumptions are already built into the logic itself. \textcite{Koellner2010SecondOrderLogic} explains in detail why full second-order logic is fragile with respect to background set-theoretical assumptions, a feature that appears to stand in tension with the traditional conception of logic as invariant across interpretations.

In my view, Feferman is correct in claiming that determinacy in these cases is ultimately accepted as a matter of faith, both in the categoricity of arithmetic and in that of set theory. Nevertheless, I maintain that the two cases differ substantially in mathematical practice. Mathematicians typically do not reason in a purely syntactic fashion, but rather \textit{within} an intended structure. Structures incorporating $\mathbb{N}$ and $\mathcal{P}(\mathbb{N}) \simeq \mathbb{R}$, within which the categoricity theorem for $\mathrm{PA}_2$ can be proved, appear ubiquitously and consistently across different areas of mathematical practice. By contrast, the full set-theoretic hierarchy rarely appears outside foundational contexts, except perhaps as an intuitive representation of the axioms of set theory themselves. Needless to say, the assumptions underlying the categoricity theorem for arithmetic are substantially weaker than those required in the corresponding case of set theory, as \textcite[141]{Ferreiros2023ConStructuralism} points out:

\begin{adjustwidth}{1cm}{1cm}
\begin{quote}
\small
The categoricity of the $\mathbb{N}$-structure can be obtained in weak second-order logic (...). Hence its fully determinate nature, so to speak. In practice, this is admitted by mathematicians and logicians who differ in their acceptance of some questionable foundational principles. The categoricity of the $V$-structure is (...) merely an ideal that guides some important research projects in advanced set theory (while it is abandoned in other projects).
\end{quote}
\end{adjustwidth}

In such cases, what we find is a commitment to a specific mathematical reality. From this reality, practitioners extract axioms that lead to the results in which they are interested. In the case of the categoricity of arithmetic, the commitment is to the structure of the natural numbers and to certain features of its subsets. These features can be formalized in various axiomatic systems, which in this case suffice to prove the desired categoricity result. In the case of set theory, by contrast, the required commitments are substantially stronger and far less common in ordinary mathematical practice.

\section{Standard models}
\label{standard-models}

The issue of non-standard models is parallel to that of categoricity, since the failure of categoricity allows the existence of non-standard models. For this reason, both first-order $\mathrm{PA}$ and first-order $\mathrm{ZFC}$ admit non-standard models. 

In the case of arithmetic, a non-standard model $M$ is characterized by the existence of a non-standard element $a \in M$, which has the feature that $a > 0_M, 1_M, 2_M, \dots$. Since its induction schema ranges only over first-order definable properties, first-order  $\mathrm{PA}$ cannot exclude non-standard elements. If $a$ is non-standard, then $a-1,a-2,\dots$ are all in $M$ and also non-standard. 

If $a$ and $b$ are such that $|a-b|$ is non-standard, then $$c = \left\lfloor \frac{a+b}{2} \right\rfloor \in M $$ is such that both $|a-c|$ and $|c-b|$ are non-standard. Defining then an equivalence relation $a\thicksim b \Leftrightarrow |a-b| \text{ is standard}$, it can be shown that each equivalence class has the order type of $\mathbb{Z}$, and the quotient order is dense without endpoints. It follows that every countable non-standard model of $\mathrm{PA}$ has order type

$$\omega + (\mathbb{Z} \cdot \eta) $$

being $\eta$ the usual order type of $\mathbb{Q}$, the set of the rational numbers. For uncountable models, $\eta$ must in general be replaced by some dense linear order $I$ without endpoints.

One feature of these models is that every non-standard model of $\mathrm{PA}$ contains, as an initial segment of its domain, a copy of the standard natural numbers. For this reason, \textcite[p.~512]{Koellner2016} states that “in the case of number theory, when you build a non-standard model, its non-standard character is immediately revealed, since the standard part is a proper substructure of the entire structure.” On the other hand, he considers the situation in set theory to be entirely analogous, describing this as a metamathematical parallelism between arithmetic and set theory:

\begin{adjustwidth}{1cm}{1cm}
\begin{quote}
\small
But the case in set theory is exactly parallel. Take a case like Cohen's method of forcing where you build a non-standard model (...). This is done within set theory and it is revealed in the first step that one is dealing with a non-standard model since it is a countable object. In the second approach one builds class size models (...) but one builds a Boolean-valued model and, once again, it is immediately revealed that it is non-standard.
\end{quote}
\end{adjustwidth}

Note, however, that while Koellner correctly identifies criteria by which a model of set theory may be recognized as non-standard, he does not address how one could positively identify a standard model.

Turning to the case of set theory in greater detail, a standard approach, as Koellner points out, is to begin with a model \(M \models \mathrm{ZFC}\) and extend it by means of an \(M\)-generic filter \(G\), obtaining a new model \(M[G]\) satisfying \(\mathrm{ZFC}\) together with some additional statement forced over \(M\). Typically, \(M\) is assumed to be both transitive and countable and, consequently, so is \(M[G]\). This is not done arbitrarily. In the first place, the Löwenheim--Skolem theorem together with Mostowski's collapse lemma allows one, in many contexts, to replace a given model by a countable transitive one without essential loss of generality. More importantly, a central result underlying forcing is the generic filter existence lemma; see, for example, Lemma IV.2.3 in \textcite[p.~246]{Kunen2011}. The standard proof of this result presupposes that the ground model \(M\) is countable from the metatheoretical perspective. Consequently, in many canonical applications of forcing---including Cohen's original construction---the models that play the central operational role are countable transitive models.

In the case of Boolean-valued models \(V^{\mathbb{B}}\), the situation is similar. In these models, sentences are assigned truth values in a complete Boolean algebra rather than the classical values \(0\) and \(1\). This is why Koellner says that such structures are ``immediately non-standard'': they do not satisfy the ordinary two-valued semantics of classical first-order logic. After choosing a \(V\)-generic ultrafilter \(G\), similarly to the previous case, the quotient \(V^{\mathbb{B}}/G\) becomes again an ordinary two-valued model satisfying \(\mathrm{ZFC}\) together with whatever additional statements are forced over the ground model. In this case, however, there is no straightforward criterion by which one could determine whether the resulting model is standard or non-standard; this depends ultimately on broader philosophical assumptions concerning the intended universe of sets.

This reveals an interesting tension concerning the notion of a ``standard model'' in set theory. In forcing practice, the structures that mathematicians most directly manipulate are not fully realized platonic universes, but countable transitive models together with their forcing extensions, or, in some cases, Boolean-valued models. In ordinary language, the term ``standard'' usually denotes something widely accepted or commonly employed. In arithmetic, although first-order theories admit non-standard models, mathematical practice has consistently privileged the standard model of the natural numbers. In set theory, by contrast, forcing practice often proceeds through surrogate structures rather than through direct engagement with the universe ordinarily regarded as standard.
This highlights a fundamental disanalogy between arithmetic and set theory. While in arithmetic, the standard model, \(\mathbb{N}\), is privileged by practice, different set-theoretic frameworks propose competing conceptions of what the universe of sets should be. Most set theorists would reject the claim that \(L\), Gödel's constructible universe, exhausts the intended universe of sets. Yet defending this position appears to require substantial additional philosophical assumptions, such as maximizing principles, which are considerably more demanding than those required to justify taking \(\mathbb{N}\) as the standard model of arithmetic. The structure \(\mathbb{N}\) is grounded in the simple iterative application of the successor operation, which appears sufficient to generate the entirety of the arithmetical universe. By contrast, the cumulative hierarchy is generated through a far more intricate network of operations---power set, replacement, transfinite recursion---which do not obviously exhaust the universe of sets.

A system of axioms such as \(\mathrm{PA}\) was originally extracted from the standard structure \(\mathbb{N}\), which then served as the basis for subsequent formal reasoning. This fits well with Conceptual Structuralism, as reflected in Thesis 3 discussed in Section \ref{conceptual-structuralism-intro}. In set theory, however, the process appears inverted: one begins with a system of axioms such as \(\mathrm{ZF}\) and then attempts to recover from them a determinate structure, namely the cumulative hierarchy of sets. But since no comparably privileged structure is available, fundamental questions often require appeal to controversial philosophical principles. Feferman appears to endorse a closely related view when he states that ``the concept of arbitrary set essential to its formulation is vague or underdetermined and there is no way to sharpen it without violating what it is supposed to be about'' (\textcite[p.~1]{Feferman2011a}).

%-------------Orey sentences---------------------------
In set theory, many have thought of independence phenomena as a reason for indeterminacy. For example, given a countable transitive model $M$ of $\mathrm{ZFC}$ we can force to obtain extensions $M[G_0]$, $M[G_0][G_1]$, $M[G_0][G_1][G_2]$,... that alternatively satisfy $CH$ and $\neg CH$, something that \textcite{Hamkins2012Multiverse} has used this fact as an argument that $CH$ should be regarded as indefinite. 

However, there is \textcite[513]{Koellner2016} a parallel situation in aritmetic. If $M \models \mathrm{PA}$ is a model of Peano Arithmetic, an end-extension of $M$ is a larger model $N \supseteq M$ such that no element of $N - M$ is inserted ``bellow'' an old element of $M$, i.e., whenever $a \in M$ and $b \in N$ satisfy $b <^N a $, then $b \in M$. If $\phi$ is an Orey sentence for $\mathrm{PA}$, then if $M$ is a model of $\mathrm{PA}$, we can end-extend it to obtain a model $N_1$ of $\mathrm{PA}$ where $\phi $ holds, and another model $N_2$ one where $\neg \phi$ does. 

First-order \(\mathrm{PA}\) is not, however, the final framework within which arithmetical truth is investigated. Goodstein's theorem and the Paris--Harrington theorem are classical examples of statements independent of \(\mathrm{PA}\), yet provable in stronger systems of arithmetic or set theory.\footnote{For Goodstein's theorem, see \textcite{KirbyParis1982}; for the Paris--Harrington theorem, see \textcite{ParisHarrington1977}.} Thus, arithmetic, no less than set theory, is subject to independence phenomena. The crucial difference, however, is that arithmetical practice has generally exhibited a stable progression toward stronger frameworks capable of resolving such statements, while preserving the privileged status of the standard structure \(\mathbb{N}\). In set theory, by contrast, the transition to stronger principles is far more philosophically contested, since no comparably privileged conception of the universe of sets appears to be available. Independence alone is not the issue; the issue is whether mathematical practice stabilizes around a privileged structure when extending axiomatic strength.

\section{So is conceptual structuralism neutral on CH?}
\label{conclusion-Koellner-clarity}

Once these points have been analyzed, we are in a better position to reassess Koellner's criticism of Feferman from a broader perspective. As discussed in Section \ref{conceptual-structuralism-intro}, Thesis 5 of Conceptual Structuralism states that full truth applies only to completely clear mathematical conceptions. According to \textcite[15]{Koellner2022}, this appeal to clarity ultimately bears most of the philosophical weight in Feferman's position:

\begin{adjustwidth}{1cm}{1cm}
\begin{quote}
\small
Finally, in the end, when all the dust settles the entire case rests on the claim that the concept of natural numbers is clear while the concept of arbitrary sets of natural numbers is not clear.
\end{quote}
\end{adjustwidth}

Koellner maintains that no purely metamathematical argument can establish an intrinsic lack of clarity in second-order arithmetic and stronger set-theoretical frameworks without threatening the determinacy of arithmetic as well. However, as argued in Sections \ref{mathematical-practice}, \ref{categoricity-second-order-logic}, and \ref{standard-models}, mathematical practice treats the corresponding structures in arithmetic and set theory in importantly different ways. In this sense, I think that stable methodological tendencies within mathematical practice provide a more robust basis for claims concerning the indeterminacy of certain mathematical structures than a freestanding notion of ``clarity'' dependent primarily on individual intuitions:

\begin{adjustwidth}{1cm}{1cm}
\begin{quote}
\small
Intuitions of clarity are not robust. Some people think that the concept of an arbitrary subset of natural numbers is about as clear a conception as one can have; others disagree. Whose intuitions are to count? \textcite[14]{Koellner2022}
\end{quote}
\end{adjustwidth}

This objection becomes even stronger if the notion of clarity itself is understood primarily in epistemological or intuition-based terms. Koellner's critique also rests on the assumption that Feferman's position depends too strongly on controversial epistemological considerations:

\begin{adjustwidth}{1cm}{1cm}
\begin{quote}
\small
The reason we should be wary of such arguments is that epistemology doesn't exactly have the greatest track record. Epistemologists have still not solved the problem of induction! But that shouldn't lead us to reject common sense claims or basic claims of physics. It should lead us to acknowledge that epistemology is a very difficult subject. Grant the limited sense in which conceptual structuralism is underdetermined. \textcite[506]{Koellner2016}
\end{quote}
\end{adjustwidth}

It seems reasonable to acknowledge that epistemology remains a highly controversial field. Consequently, any strong philosophical argument concerning the definiteness of CH would remain equally controversial insofar as it depended too heavily on disputed epistemological assumptions. This is precisely why the appeal to mathematical practice developed throughout this paper is important. The asymmetries between arithmetic and higher set theory need not rest exclusively on controversial intuitions concerning clarity or epistemological access, but can instead be grounded in the comparatively stable ways these structures function within ordinary mathematical and scientific practice.

Due to how concisely Feferman addressed this topic, as explained in Section \ref{conceptual-structuralism-intro}, it remains partly a matter of interpretation how strongly his position depends on epistemological considerations, and how much weight it instead assigns to methodological considerations grounded in mathematical practice. For example, Feferman acknowledged that set theory possesses an intuitive representation in the cumulative hierarchy, while nevertheless refusing to commit himself either to the existence of the entire set-theoretic universe or to a fully platonistic conception of mathematics:

\begin{adjustwidth}{1cm}{1cm}
\begin{quote}
\small
ZF has an intuitive model in the transfinite iteration of the power set operation taken cumulatively. This has nothing to do with a belief in a platonic reality whose members include the natural numbers and arbitrary sets of natural numbers, and so on. On the contrary, I disbelieve in such entities. \textcite[72]{Feferman2000rationale}
\end{quote}
\end{adjustwidth}

Perhaps Feferman is here distinguishing between the legitimacy of employing the cumulative hierarchy as a mathematical or methodological framework and commitment to the existence of a unique, fully determinate universe of sets. In this sense, one may coherently reason within the framework provided by ZF without thereby accepting a platonistic interpretation of its ontology. Feferman appears to express a similar idea in a later discussion of CH:

\begin{adjustwidth}{1cm}{1cm}
\begin{quote}
\small
There is no problem to put oneself in the mental frame of mind of “this is what the cumulative hierarchy looks like”, for which one can see that such and such propositions including the axioms of ZFC are (more or less) obviously true. I have taught set theory many times and have presented it in terms of this ideal-world picture with only the caveat that this is what things are supposed to be like in that world, rather than to assert that that's the way the world actually is. \textcite[19]{Feferman2011a}
\end{quote}
\end{adjustwidth}

Feferman's insistence on clarity can also be interpreted as a modern continuation of the predicative tradition, although he did not always explicitly identify with it:

\begin{adjustwidth}{1cm}{1cm}
\begin{quote}
\small
Though I am not now nor never have been a predicativist, I have to admit to being a sympathizer since I am an avowed anti-platonist, at least insofar as set theory is concerned, and I grant the natural numbers a position of primacy in our mathematical thought. \textcite[313]{Feferman2004a}
\end{quote}
\end{adjustwidth}

Feferman does not explicitly formulate his position in these terms, but his concerns about arbitrary totalities may be interpreted as belonging to the same broader philosophical tradition. Let us recall the liar paradox, expressed by the sentence ``This statement is false''. \textcite[224]{Russell1908} rephrased it as follows: ``It is not true for all propositions $p$ that if I affirm $p$, $p$ is true''. He then concluded that the source of the paradox lay in quantification over all propositions, so that the notion of an arbitrary proposition appeared deeply problematic. A similar line of reasoning was employed in discussions of the barber paradox, which arises from attempting to define the \emph{naïve} set $\{x : x \notin x\}$. The unrestricted notion of an arbitrary naïve set was regarded as insufficiently clear or determinate, and this lack of clarity was taken to underlie the resulting contradiction. Intuitionists extended this critique even further by arguing that such paradoxes arise from accepting logical structures that lack a counterpart in intuition, such as the law of excluded middle.

Perhaps Feferman's attitude toward CH can be interpreted in a broadly similar spirit. On this reading, the issue would not be that the concept of an arbitrary subset of the continuum is straightforwardly contradictory, but rather that it lacks the kind of conceptual grounding required for full determinacy. In this sense, Feferman's concern may resemble earlier predicative and intuitionistic worries about unrestricted totalities: namely, that certain forms of quantification extend beyond what can be regarded as sufficiently clear or conceptually secured. If this interpretation is correct, then Feferman's position on CH should not be understood primarily as a claim about inconsistency, but rather as skepticism concerning whether the relevant set-theoretical totalities are determinate in the same way as ordinary arithmetical structures.

\section{Conclusion}
\label{conclusion-chapter1}

The Feferman--Koellner dispute concerning CH forms part of a long tradition of philosophical and mathematical controversies generated by the Continuum Hypothesis since Cantor first posed it. Both Feferman and Koellner have made precise and insightful contributions to these controversies. In my view, however, Koellner's criticism underestimates the extent to which Feferman's conceptual structuralism can be interpreted through the lens of mathematical practice.

In this paper, I have argued that Feferman's conceptual structuralism can be interpreted in a way that withstands some of Koellner's strongest criticisms. While Koellner is correct that metamathematical considerations alone do not establish the indefiniteness of the Continuum Hypothesis, this does not undermine Feferman's broader philosophical position. The crucial point is that questions concerning mathematical determinacy cannot be settled exclusively through categoricity results, formal semantics, or logical reconstruction. They must also take into account the role mathematical structures play within scientific and mathematical practice.

From this perspective, arithmetic and higher set theory occupy importantly different roles. The structure of the natural numbers is deeply integrated into ordinary mathematical reasoning, exhibits remarkable stability across foundational frameworks, and is closely connected to basic cognitive and operational practices. By contrast, the higher set-theoretic universe appears far more framework-dependent, less directly connected to ordinary mathematical activity, and more strongly shaped by optional foundational commitments. These asymmetries do not by themselves prove that CH is indefinite, but they help explain why its determinacy remains philosophically contested in a way that ordinary arithmetical truth does not.

The resulting picture is neither straightforward platonism nor simple conventionalism. Rather, mathematical determinacy emerges gradually from the interaction between formal systems, conceptual clarity, and the stabilizing role of mathematical practice itself. Interpreted in this way, conceptual structuralism offers not merely a skeptical attitude toward CH, but a broader philosophical framework for understanding why some mathematical structures come to be treated as fixed and objective while others remain open to competing interpretations.

%-----------Segundo capítulo--------
\chapter{The idea of \textit{quasi-empirical structuralist} and its implications to the Continuum Hypothesis}
\label{chapter2}

%=========================
\section{Introduction}
\label{Intro-chapter-2}

Toward the end of the 19th century, Georg Cantor developed a theory of cardinality for infinite sets. Among his key results was the proof that the cardinality of the continuum strictly exceeds that of the natural numbers. This prompted him to ask whether any set could exist with a cardinality strictly between the two --- a conjecture that came to be known as the \textit{Continuum Hypothesis} (CH). The argument generalizes naturally: for any set $A$, one can show that the cardinality of its power set $\mathcal{P}(A)$ is strictly larger than that of $A$ itself. This observation gives rise to the \textit{Generalized Continuum Hypothesis} (GCH), which asserts that no set of intermediate cardinality can exist between $A$ and $\mathcal{P}(A)$. 

The work of \textcite{Godel1940} and \textcite{Cohen1966} showed it to be independent of the ZFC axioms of set theory, and therefore impossible to decide using the tools usually available to mathematicians in practice. This has led some authors to question whether the question admits a unique and meaningful answer at all. Among the most prominent is Solomon Feferman, who has argued (\citeyear{Feferman2011a}) that $\mathrm{CH}$ is an indefinite mathematical statement---the mathematical structures in which it could be evaluated being, in his view, not clear and intrinsically vague.

In this paper I would like to face this controversy from a different angle. In my view, one of the most crucial and delicate features of any meaningful mathematical statement is whether it amounts to more than a mere logical tautology. Almost any mathematical result can of course be read as a true statement within a given formal system---but the real significance of mathematics lies in whether its symbols connect to something beyond that system, whether they carry meaning that is not exhausted by the rules governing them. This is a central concern of the \textit{philosophy of mathematical practice}, a growing field that takes seriously the gap between formal proof and mathematical understanding. 

This paper is guided by two complementary philosophical commitments. The first is the quasi-empiricist philosophy of mathematics developed by \textcite{Lakatos1976}, according to which mathematical knowledge grows through an iterative process of conjecture, refutation, and revision driven by practice itself—not by the top-down imposition of formal systems. In this sense, mathematical fruitfulness becomes a central guide to the development of mathematical knowledge: concepts, structures, and methods earn their place not merely through formal consistency, but through their capacity to solve problems, generate new lines of inquiry, and contribute to mathematical understanding. Underlying this view is a broader claim about mathematical meaning: that the significance of a mathematical statement cannot be read off from its position within a formal system alone, but depends on the uses to which it is put and the problems it is called upon to solve. This idea finds a striking expression in Wittgenstein:

\begin{adjustwidth}{1cm}{1cm}
\small
\begin{quote}
It is essential to mathematics that its signs are also employed in mufti (...). It is the use outside mathematics, and so the meaning of the signs, that makes the sign-game into mathematics. \parencite[V §2]{Wittgenstein1978}\footnote{Wittgenstein uses the phrase “employed in mufti” to mean that mathematical signs are used outside their formal mathematical setting—in ordinary contexts such as science, measurement, engineering, or everyday reasoning. In these contexts, the symbols appear without their \textit{formal mathematical uniform}, that is, outside the strictly formal system.}
\end{quote}
\end{adjustwidth}

%------------------------------------
The second commitment is a structuralist approach to mathematical discourse and methodology. On this view, mathematical statements are formulated within structures whose relevant features are determined by the relations among their constituent objects rather than by the intrinsic nature of those objects. Whether one interprets such structures realistically, as \textcite{Shapiro1997} does, or more modestly as patterns in Resnik's (\citeyear{Resnik1981}) sense, mathematical meaning is inseparable from the structures within which mathematical discourse takes place.

This perspective motivates a distinction that will play a central role in the argument of this paper. A mathematical structure is adopted when mathematical practice privileges a particular realization as capturing the intended object of study, so that alternative realizations are treated merely as representations of the same structure. By contrast, a pattern is adopted when different realizations are treated as equally legitimate mathematical possibilities and no particular realization is singled out as the intended one. The distinction is methodological rather than ontological: it concerns how mathematical practice relates to structures and axioms, rather than what ultimately exists.

This last point bears directly on the debate surrounding the Continuum Hypothesis. It is of course true that $\mathrm{CH}$ has a precise and well-defined meaning within the language of set theory. But well-definedness within a logical language does not, in general, suffice for genuine mathematical significance---if it did, even the most elaborate and idle tautologies would qualify as mathematical claims. Set theorists routinely investigate a wide variety of such structures, many of which are exceptionally rich and well understood. The philosophical question is rather whether mathematical practice provides grounds for privileging one of these structures as the intended universe of sets. In particular, the Continuum Hypothesis is true in Gödel's constructible universe $L$ and false in many forcing extensions. The issue is therefore not whether $\mathrm{CH}$ can be evaluated in a determinate structure, but whether mathematical practice provides sufficient grounds for privileging one such structure as the intended universe of sets.

The role of fruitfulness in the present account is indirect. Mathematical practice does not confer truth values on statements; rather, it guides the adoption of mathematical structures. Once a structure has been adopted as sufficiently determinate, statements formulated within it inherit determinate truth conditions. The central question is therefore not whether $\mathrm{CH}$ has a truth value in some set-theoretic universe, but whether contemporary mathematics has adopted a privileged set-theoretic structure in the same way that it has adopted the standard structures of arithmetic and analysis. My claim is that, at present, there is insufficient reason to think that it has. The concern, then, is not that $\mathrm{CH}$ lacks a formal interpretation, but that formal interpretability alone is insufficient for mathematical significance unless mathematical practice provides grounds for privileging one such structure.

The remainder of the paper is organized as follows. Section \ref{structuralims-ontology-epistemology} develops the structuralist framework underlying the present account and introduces the distinction between the adoption of mathematical structures and the adoption of axiom systems. Section \ref{quasi-emp-structuralism} turns to the philosophy of mathematical practice, drawing on Lakatos and Mac Lane to formulate the quasi-empirical criteria of mathematical fruitfulness that guide the adoption of structures. Section \ref{history-section} illustrates these ideas through a series of historical case studies, including the development of complex numbers, function spaces, and infinitesimals. Section \ref{natural-reals} argues that the standard structures of arithmetic and analysis have achieved a privileged status within mathematical practice, thereby supporting determinate truth conditions for statements formulated within them. Section \ref{set-theory} contrasts this situation with that of contemporary set theory, arguing that mathematical practice has largely adopted set-theoretic methods and axioms without thereby privileging a unique set-theoretic universe. Section \ref{section-CH-case} applies the preceding analysis to $\mathrm{CH}$ and examines whether either contemporary axiomatic developments or current mathematical practice provide grounds for treating $\mathrm{CH}$ as determinate in the relevant sense.

\section{The structuralists' approach: ontology, epistemology and categoricity}
\label{structuralims-ontology-epistemology}
 
An important philosophical approach that distances itself from formalism while engaging seriously with both epistemological and ontological questions is mathematical structuralism, developed most prominently by Shapiro and Resnik. As will become clear, structuralist approaches---including Feferman's conceptual structuralism---bear directly on debates about the determinacy of set-theoretic statements such as CH, particularly the question of whether there is a single, fixed set-theoretic universe.

One of the most paradigmatic developments within this philosophical movement is that of \textcite{Shapiro1997}. In this work, he defends a conception known as \textit{ante rem structuralism}, according to which mathematical structures exist independently of—and prior to—the systems of objects that instantiate them. On this view, what matters is not which particular objects occupy positions within a given structure, but rather the relations that define the structure and the ways in which structures relate to and extend one another.

Shapiro remains nevertheless open to alternative forms of structuralism, acknowledging them as broadly equivalent in practice even if less natural with respect to the actual development of mathematics. This includes Resnik's pattern-based (\citeyear{Resnik1981}) or non-ontological structuralism (\citeyear{Resnik2018NonOntological}). Shapiro himself is explicit on this point:

\begin{adjustwidth}{1cm}{1cm}
\small
\begin{quote}
I believe that the \textit{ante rem} option is the most perspicuous and least artificial of the three [eliminative, modal eliminative and \textit{ante rem} structuralism]. It comes closest to capturing how mathematical theories are conceived. Nevertheless, I do not mean to rule out the other options. Indeed, it follows from the thesis of structuralism that, in a sense, all three options are equivalent. \parencite[p. 90]{Shapiro1997}
\end{quote}
\end{adjustwidth}

%--------Important upshot of structuralism: in a really weak form--------------
The most important upshot of ante rem structuralism, for the purposes of this paper, is that structures are the fundamental primitives of mathematics, and that mathematicians characteristically reason in structural terms. Chapters 5 and 6 of \textcite{Shapiro1997} illustrate how deeply this perspective is embedded in the history of mathematics. This is evident, for instance, in \textcite{Dedekind1888}, as well as in \textcite{Bourbaki1950}, where “structure” is conceived as a collection endowed with one or more relations satisfying a specified set of axioms. The philosophical consequence is a broadly holistic constraint on mathematical meaning: a statement can only be considered determinate or meaningful relative to the structures in which it is defined.

Considerations around structuralism naturally raise the question of whether mathematical structures exist prior to objects, but also introduce a further “chicken or egg” problem: are structures prior to the axioms that characterize them, or are they instead determined through axiomatic specification? \textit{Ante rem} structuralism suggests that axioms merely describe independently existing structures, although the position is less clear on this issue than on the question of object-structure priority.

It is at this point where the epistemological problems regarding structures come into play. According to \textcite{Shapiro1997}, and much in line with \textcite{Resnik1981}, the fundamental way of coming to know a mathematical structure is through pattern recognition. The clearest case is the natural number structure \parencite[pp. 118-119]{Shapiro1997}, where the relevant pattern is the indefinitely iterable act of adding one element to a finite sequence.

At this point we see the sharp line between pattern recognition and axiom systems. When more abstract mathematical structures, such as the set theoretical hierarchy, come into play, then the patterns leading to them may be more complex. At this point, the pattern involved leads to fundamentally an axiom system, such as $\mathrm{ZFC}$, and the structure is apprehended not directly but through what Shapiro calls implicit definition:

\begin{adjustwidth}{1cm}{1cm}
\begin{quote}
  \small
One way to understand a particular structure is through a direct description of it (...). Eventually, a properly prepared listener will understand that it is the structure itself, and not any particular instance of it, that is being described (...). We have here an instance of implicit definition, a technique familiar in mathematical logic. \parencite[p. 130]{Shapiro1997}
\end{quote}
\end{adjustwidth}

According to Shapiro, thus, relevant structures in mathematical discourse and axioms describing them need to stay at the same level, which in logical terms we understand as having structures determined by categorical axiom systems. However, crucial axiom systems for the practice of mathematics, such as $\mathrm{PA}$ for arithmetic and $\mathrm{ZF}$ for set theory, are not categorical in first order logic.\footnote{As categoricity of a theory implies completeness, any theory affected by Gödel's incompleteness theorem won't be categorical.} Accordingly, \textcite{shapiro1991foundations} argued that second order logic, where these theories are categorical (quasicategorical in the case of set theory), provides a suitable foundation for mathematics. In chapter 8 of this book, Shapiro defends that, in pre-formal practice, mathematicians operate with concepts such as ``all subsets'', ``all natural numbers'', that fixes a unique interpretation for the axioms which enable them to distingish stardard from non standard models.

This perspective, however, is subject to some important criticism. \textcite{Vaananen2015}---much in line with \textcite{Quine1941}---argues that second-order logic depends on what he calls large domain assumptions, which come quite close to the meaning of the axioms of set theory. \textcite{Shapiro2012} even accepts this, but holds that it does not undermine his case for second order logic as a natural foundation of mathematics. They agree, for that reason, that second order logic and first order set theory can be threated as logically equivalent. However, from the perspective of the practice of mathematics, the latter is no doubt more prominent in the various foundationalist approaches to mathematics. 

Arguments that Shapiro gives related to mathematical practice are also problematic. For example he states (\citeyear{shapiro1991foundations}, p. 99) that important tool in mathematical practice, which is the concept of a minimal closure, cannot be formalized without appealing to second order logic: ``There is a straightforward argument, relying on the compactness theorem, that no collection of first-order formulas can successfully characterize any non-trivial minimal closure. This is evidence that first-order languages are inadequate.'' However, applying this argument to mathematical practice is not that straightforward, since in general practicioners tend to use minimal closures that are defined within a previously given structure. In a given structure, most closures are definable by formulas with parameters or describable as the intersection of definable sets. Shapiro's compactness argument establishes a genuine limitation of first-order logic in the abstract, but it does not map cleanly onto the way closures are actually used.\footnote{The algebraic closure of $\mathbb{Q}$ provides a familiar example, working within the structure $\mathbb{C}$ as given.}

%-------------Distinction: open-ended vs determined structures----------------
In this discussion---the priority of axioms and patterns, or that of structures over them---not all the structuralists agree Shapiro. \textcite[Chapter 7]{Resnik1997} defends that a mathematical theory, given by axioms, merely describes a structural pattern, without necessarily determining a unique and privileged structure that models that theory.

A clear and paradigmatic case of this approach in mathematical practice is category theory. With a remarkably concise set of axioms---objects, morphisms, composition, and identity---the theory applies across a wide variety of mathematical domains, including groups, topological spaces, sets, relations, partially ordered sets, and differentiable manifolds. What matters in practice, however, are the particular collections that instantiate these axioms, rather than whether there exists a fundamental and privileged structure that realises them. This is precisely the spirit in which \textcite{maclane1998categories} develops the theory: category theory is presented as a unifying language for mathematical practice, one whose power lies in its ability to reveal structural analogies across disparate areas, not in grounding mathematics in a single canonical structure. MacLane is notably cautious about foundational overreach---he treats the axioms of category theory as a flexible and powerful descriptive tool rather than as the implicit definition of a unique intended model.\footnote{Similarly, \textcite{Feferman2013Category} argues that the strength of category theory lies precisely in what MacLane emphasised---its role as a structural lingua franca across mathematical domains---rather than in any foundational pretension. He is also critical of the notion of a ``category of all categories'': beyond the coherence problems such a universal structure faces, he argues that explaining what a category is already presupposes a prior notion of collection and operation, so category theory cannot serve as a logically autonomous foundation.}

The category theory case thus illustrates, in a particularly vivid way, the broader tension between the two modes of structural aprehension identified above. Neither of these two approaches---the priority of structures, or the priority of axioms and patterns---is unanimously adopted in mathematical practice, and I do not think one should be universally privileged over the other. Rather, the two capture genuinely different modes of mathematical commitment, and the distinction between them is philosophically consequential. 

The disagreement between Shapiro and Resnik is not merely ontological. It also reflects different ways in which mathematical practice relates structures to the axioms used to describe them. In some cases, mathematicians appear to begin with a structure that is already understood, at least informally, and subsequently formulate axioms intended to capture its essential features. In other cases, they begin with an axiom system and investigate the range of structures satisfying it, without taking any particular model to be privileged. The distinction between these two situations is not exhausted by questions of ontology or epistemology; it concerns the role that structures and axioms play within mathematical practice itself. To mark this distinction, I introduce the notions of adoption of a mathematical structure and adoption of an axiom system.

When mathematicians come to understand a mathematical structure prior to its formal axiomatization—formulating axioms only later as descriptions of something already grasped in practice—I shall say that they have \textit{adopted a mathematical structure}. This usage is broadly consonant with the spirit of \textit{ante rem} structuralism, according to which structures are not created by their axioms but are instead captured by them. By contrast, when mathematicians begin from an axiom system before any particular structure has been singled out as the intended object of study—a methodological tendency that has become increasingly common since the formalist turn—I shall say that they have \textit{adopted a pattern}. In such cases, priority is given to the axiomatic specification rather than to any particular realization of it. When the axioms are non-categorical, this adoption does not determine a unique structure: many non-isomorphic models satisfy the same axioms, and mathematical practice provides no principled grounds for privileging one of them as the intended realization. By Gödel's completeness theorem, every consistent first-order axiom system has a model; assuming mathematicians commit only to systems they regard as consistent, adopting such a pattern amounts to treating the entire class of models satisfying the axioms as legitimate mathematical possibilities. I shall refer to this class, throughout the remainder of the paper, as an \textit{open-ended structure}.

The distinction, as pointed out in the \nameref{Intro-chapter-2}, bears directly on the question of CH. When mathematicians adopt a mathematical structure, statements formulated in the corresponding language are treated, at least on paper, as possessing determinate truth values, even when those values cannot be established by currently available proof methods. By contrast, adopting a non-categorical axiom system amounts only to adopting a pattern that may be realised by many non-isomorphic structures. In that case, semantic determinacy is not secured by the axioms alone: a statement may receive different truth values in different models. Questions such as $\mathrm{CH}$ therefore remain open not merely epistemically but also structurally, insofar as no privileged structure is available to fix a unique answer.

%=========================
\section{Stepping away from formalism: mathematical practice and metodology}
\label{quasi-emp-structuralism}

As we have seen, structuralism provides important ontological and epistemological insights into mathematics, particularly through the notions of adoption of mathematical structures and axiom systems. Our concern here, however, is not merely the nature of such commitments, but the reasons mathematicians adopt particular structures or axioms in actual mathematical practice.

Philosophical approaches to mathematics have evolved dramatically over the course of the twentieth century. At the beggining of this century, both the mathematical world—strongly influenced by Hilbert’s formalism—and several philosophical approaches to mathematics, such as logical positivism, located the essence of mathematics in formal and logical methods. From the perspective of formalism, mathematical statements were understood as deductions from a given set of axioms, and any consistent theory formulated in a logical language was regarded as sufficient to define a corresponding piece of mathematical reality—see, for instance, \textcite{hilbert1971geometry}. Within logical positivism, as represented for example by \textcite{carnap1937syntax}, mathematical sentences were regarded as analytic, in the sense that they say nothing about the empirical world and function instead as tools for clarifying the structure of scientific language.

In the latter half of the twentieth century, new philosophical conceptions of mathematics emerged that emphasized practical applicability, heuristic reasoning, and the essentially informal character of everyday mathematical work. As pointed out in the \nameref{Intro-chapter-2}, this shift marked the advent of the so-called \textit{philosophy of mathematical practice}. One of its earliest and most influential proponents was Imre Lakatos, who, in his work \textit{Proofs and Refutations}, strongly criticized the formalist approach of earlier schools:

\begin{adjustwidth}{1cm}{1cm}
\small
\begin{quote}
Since, [according to logical positivism], informal mathematics is neither `tautological' nor empirical, it must be meaningless, sheer nonsense. The dogmas of logical positivism have been detrimental to the \textit{history and philosophy of mathematics}. \parencite[p. 2]{Lakatos1976}
\end{quote}
\end{adjustwidth}

Lakatos’ alternative approach is known as quasi-empiricism, as he argued that mathematics develops in a manner analogous to the empirical sciences: one begins with a conjecture and subsequently seeks a proof. Although such a proof may initially appear convincing, it can generally be decomposed into lemmas and implicit assumptions, which may themselves be vulnerable to refutation. When such refutations arise, the original conjecture is not simply discarded but rather refined—whether by restricting its scope, revising underlying definitions, or modifying auxiliary assumptions.\footnote{This idea is extensively developed through the analysis of Euler's formula, which states that for convex polyhedra with $V$ vertices, $E$ edges, and $F$ faces, $V - E + F = 2$. See \textcite[pp. 10–41]{Lakatos1976}.}

According to Lakatos, this perspective remains meaningful even in the context of the increasing rigor of nineteenth-century mathematics, particularly with the formalization of analysis:

\begin{adjustwidth}{1cm}{1cm}
\small
\begin{quote}
In the early nineteenth century the flood of counterexamples brought confusion. Since proofs were crystal-clear, refutations had to be miraculous freaks, to be completely segregated from the indubitable proofs. Cauchy’s revolution of rigour rested on the heuristic innovation that the mathematician should not stop at the proof: he should go on and find out what he has proved by enumerating the exceptions, or rather by stating a safe domain where the proof is valid.\parencite[p. 55]{Lakatos1976}
\end{quote}
\end{adjustwidth}

From the perspective developed here, Lakatos draws a clear distinction between mathematical reasoning and purely formal calculation. Rather than grounding mathematics in formal derivations alone, he emphasizes its problem-solving power: definitions, entities, and assumptions are justified insofar as they contribute to the advancement of inquiry. Another central feature of his view is the iterative character of mathematical reasoning. The emergence of a refutation or the identification of a flaw does not terminate the process; instead, it prompts the introduction of revised assumptions, refined definitions, or new intermediate results, which may in turn generate further tensions and revisions. This iterative character is central to Lakatos’ account, according to which mathematics advances not through the accumulation of final proofs, but through a dynamic process of conjecture and criticism. As he puts it (\citeyear{Lakatos1976}, p.5), “mathematics [\dots] does not grow through a monotonous increase of the number of indubitably established theorems but through the incessant improvement of guesses by speculation and criticism, by the logic of proofs and refutations”.

%--------Inserting MacLane-------------------
Lakatos' approach illuminates not just how mathematics is justified, but how it develops---through problems, failures, and the iterative refinement of concepts. This raises a question central to the present paper: how do mathematical structures and axiom systems actually come into being, and why are they introduced when they are? This is a question that, in a similar vein to Lakatos, much oriented to practice, \textcite[36]{maclane1986mathematics} further explores:

\begin{adjustwidth}{1cm}{1cm}
\small
\begin{quote}
The genesis of the more complex mathematical structures tends to take place within Mathematics itself. [There] are a variety of processes which may generate new ideas and new notions. We list a few of these processes in tentative form for further refinement after our more detailed studies. \parencite[p. 36]{maclane1986mathematics}
\end{quote}
\end{adjustwidth}

MacLane's list includes \textit{conundrums} (the pressure exerted by hard unsolved problems), completion, invariance, common structure (the recognition of structural analogies across different areas), generalisation, abstraction, axiomatisation, and analysis of proof. What is striking about this catalogue is how closely it mirrors the Lakatosian picture: mathematical structures do not arrive fully formed, imposed from above by logical fiat, but are drawn out by the demands of practice. MacLane makes this developmental dimension explicit later in the same work, arguing that mathematics emerges from human experience and scientific practice and grows through successive cycles of formalisation and generalisation:

\begin{adjustwidth}{1cm}{1cm}
\small
\begin{quote}
Mathematics, in our description, rests on ideas and problems arising from human experience and scientific phenomena and consists in many successive and interconnected steps in formalizing and generalizing these inputs. \parencite[449]{maclane1986mathematics}
\end{quote}
\end{adjustwidth}

%--------------Externally applicable structures------------------------
The reference to human experience and science is crucial here. MacLane's account of how mathematics develops has the shape of a mathematical induction. There is a base case: a set of notions directly grounded in practical and scientific engagement with the physical world---measurement, counting, spatial reasoning, the mathematics demanded by physical problems. From this foundation, the iterative processes MacLane catalogues---generalisation, abstraction, the recognition of common structure, the analysis of proof---allow the totality of mathematical knowledge to grow, each step building on what came before while progressively loosening its ties to the physical ground. What results is a body of mathematics that is largely self-sustaining and internally driven, yet never fully detached from the base that gave rise to it.

This duplicity---between mathematics as an autonomous discipline pursued for its own sake and mathematics as something that cannot afford to lose contact with the external world---is captured with remarkable precision by \textcite[pp. 332--337]{Poincare1897rapports}:

\begin{adjustwidth}{1cm}{1cm}
\small
\begin{quote}
And above all, those devoted to mathematics find in it pleasures analogous to those given by painting and music (\dots). That is why I do not hesitate to say that mathematics deserves to be cultivated for its own sake, and that theories which cannot be applied to physics should be esteemed just as highly as the others (...). But the pure mathematician who would forget the existence of the external world would be like a painter who knew how to harmoniously combine colors and forms, but lacked models. His creative power would soon run dry. 
\end{quote}
\end{adjustwidth}

How might these insights be integrated into the perspective developed here, that assumes a structuralist approach to mathematical discourse and metodology? I propose to do so through the notion of $\mathcal{S}$, the---possibly open-ended---collection of mathematical structures to which mathematical practice is committed. Following the induction analogy introduced above, $\mathcal{S}$ must have a base case: a set of structures whose appearance in practice is grounded directly in their indispensability, or at least their deep entrenchment, in scientific and practical engagement with the physical world. This does not require accepting the indispensability argument of \textcite{quine1960} and \textcite{putnam1971} in its full ontological strength---the claim is not that mathematical objects exist because science needs them. The weaker claim defended here is simply that structures which have proven most successful in modelling physical phenomena, and which are thereby thoroughly embedded in scientific practice, should belong to any reasonable conception of mathematics. A mathematical structure satisfying these conditions will be called an \textit{externally applicable structure}.

%-----Internally fruitful structures-------
Beyond the base case, $\mathcal{S}$ is closed under the following principles. The first three criteria characterize mathematical fruitfulness and capture, in schematic form, central features of the approaches of Mac Lane and Lakatos. The fourth concerns a distinct notion, namely the adoption of a structure or axiom system within mathematical practice.:

\begin{enumerate}

    \item \label{condition1} \textbf{Problem-solving power:} $X$ provides solutions to previously existing, well-defined, and relevant problems formulated within structures in $S$.

    \item \label{condition2} \textbf{Explanatory unification and simplification:} $X$ reveals common structural principles underlying phenomena previously treated separately, often through abstraction and axiomatization, thereby simplifying existing theories and enabling the transfer of results and methods across structures in $S$.

    \item \label{condition3} \textbf{Generative capacity:} $X$ gives rise to new and fruitful lines of inquiry within other structures in $S$.
    
    \item \label{condition4} \textbf{Stability of reliance:} the structure or axiom system $X$ becomes part of the shared background assumptions on which substantial areas of mathematics routinely depend.
    
\end{enumerate}

Criteria \ref{condition1}, \ref{condition2}, and \ref{condition3} are best understood as criteria of mathematical fruitfulness rather than criteria of adoption. They explain why a structure or axiom system becomes a serious object of mathematical investigation, but not why it becomes part of the background framework of mathematical practice itself. Adoption is a stronger notion, referring to those structures and axiom systems upon which mathematicians routinely rely in their ordinary work.

The distinction is important. A conjecture may possess considerable explanatory or problem-solving power and nevertheless remain a hypothesis rather than becoming part of the accepted background of mathematical reasoning. Otherwise, any sufficiently fruitful but unresolved conjecture would eventually have to be adopted simply because of its mathematical utility, something that mathematical practice clearly does not tend toward. For this reason condition \ref{condition4} should be included.

The perspective developed here may therefore be understood as a form of \textit{quasi-empirical structuralism}: a structuralist conception of mathematics in which the development of mathematical structures is closely connected to their explanatory, problem-solving, and generative role within both mathematical and scientific practice. I think this conception is also faithful to the \textcite{Wittgenstein1978} passage mentioned in the \nameref{Intro-chapter-2}, insofar as it requires mathematical structures to be applicable to problems, phenomena, and areas of inquiry beyond their own immediate definition, and ultimately beyond mathematics itself.

What quasi-empirical structuralism defends, then, is a kind of methodological well-foundedness for mathematical structures. If $S_1$ and $S_2$ are mathematical structures,\footnote{I insist that such structures may be open-ended, meaning—see the end of section \ref{structuralims-ontology-epistemology}—that they may consist simply of collections of axioms accepted in practice.} I write 

$$S_1 \leadsto S_2 $$

to indicate that $S_2 \in \mathcal{S}$ because it fulfulls condition \ref{condition4} and one among \ref{condition1}, \ref{condition2} or \ref{condition3} with respect to $S_1 \in \mathcal{S}$. In this case I would say, in a way that is natural regarding those conditions, that $S_2$ is \textit{internally fruitfull} with respect to $S_1$.\footnote{Externally applicable structures fulfill this conditions with respect to real-life processes and conceptions of our sensorimotor cognition.} Then for any structure $S \in \mathcal{S}$, there is a sequence 

$$ S_0 \leadsto ... \leadsto S_n = S $$

where $S_0$ is an externally applicable structure. What concerns us here, then, is whether the structures appearing in such a sequence, in the particular case where $S$ is involved in the formulation of $CH$, are open-ended or definite. More precisely, the question is whether the fruitfulness of those structures in mathematical practice derives primarily from the adoption of definite mathematical objects, or instead from the continued development of inferential and axiomatic patterns that remain open-ended while guiding mathematical activity.


\section{Some historical examples}
\label{history-section}

In this section I examine several historical examples drawn from mathematical practice. The main objective is to argue that the iterative component of quasi-empirical structuralism has been present at various moments in the historical development of mathematics. One of the most paradigmatic cases is the development of complex numbers from the sixteenth to the nineteenth century. Initially regarded as obscure, metaphysical, or ill-defined, complex numbers eventually became one of the central structures of modern mathematics. Before turning to the implications of this example, I outline the reasons for this transformation.

%---------------Complex numbers history--------------------------
The early stages of this development can be traced back to the sixteenth century. Although the idea of negative square roots had already been mentioned before, it was Rafael Bombelli, in his 1572 work L’Algebra, the first to systematically use of what we now call complex numbers.\footnote{Most of the events mentioned here are important facts in history of mathematics, where I am following texts on the topic such as \textcite{Nahin2010} and \textcite{Christen2018}.} Bombelli was working on the solution of the reduced cubic equation, an equation of the form $x^3 + px + q = 0$. Earlier, Cardano had found the following formula to solve such equations:

\begin{equation}
x = \sqrt[3]{-\frac{q}{2} + \sqrt{\left(\frac{q}{2}\right)^2 + \left(\frac{p}{3}\right)^3}} 
+ \sqrt[3]{-\frac{q}{2} - \sqrt{\left(\frac{q}{2}\right)^2 + \left(\frac{p}{3}\right)^3}}
\end{equation} 

Cardano’s difficulty lay in interpreting this formula when negative square roots were involved. In fact, he concluded that the formula was not solvable “in general.” Bombelli decided to step a little bit further, introducing the symbol $pdm$ to represent $\sqrt{-1}$, thereby providing a clear precursor to the imaginary unit.\footnote{This stands literally for \textit{piú di meno}, which mean “plus or minus” in Italian—see \textcite{Christen2018} for more details. He arrived at many of the operational properties of complex numbers, such as $pdm \cdot pdm = -1$.} In this way, he was able to obtain new real roots of certain equations of the given form, which Cardano had been unable to.

Progress in this area accelerated in the following years. While Descartes showed little interest in such numbers, referring to them pejoratively as “imaginary,” Leibniz used them in a manner similar to Bombelli, as intermediate steps in solving equations. Euler was probably the first to use them systematically: he introduced the notation $i$ for the imaginary unit and employed it in many formulas. Perhaps the most well known is Euler's formula, which states that for any real phase $\alpha$, $e^{i\alpha} = \cos \alpha + i \sin \alpha$. Another notable application is the factorization of fourth-degree polynomials with no real roots (such as $x^4 + 1$). Later, Gauss introduced the complex plane, standardized the notation, and firmly established the acceptance of complex numbers, using them in many important results, most notably in the fundamental theorem of algebra. Since then, complex numbers have become one of the cornerstones of mathematics, and the set of complex numbers is now regarded as a mathematical structure in its own right.

%------------Reasons for shift------------------
What were the reasons for such a shift? In my view, complex numbers constituted the simplest algebraically closed extension of the real line while satisfying the three criteria of mathematical fruitfulness introduced in Section \ref{quasi-emp-structuralism}. First, they exhibited significant problem-solving power, as seen in the factorization and solution of higher-degree polynomial equations. Second, they provided explanatory unification: the Fundamental Theorem of Algebra and the unification of trigonometric and exponential functions through Euler's formula are particularly striking examples. Third, they generated new and fruitful lines of inquiry through their deep connections with geometry, analysis, algebra, number theory, and many other areas of mathematics. In this sense, the emergence of complex numbers provides a paradigmatic instance of the relation $S_1 \leadsto S_2$ introduced above. At least in their earliest applications to algebraic equations, complex numbers (endowed with their algebraic operations) may be regarded as internally fruitful with respect to the real number structure:

$$ (\mathbb{R}, +, \cdot) \leadsto (\mathbb{C}, +, \cdot) $$

%-------------------Arbitrary functions & Wave equation--------------------

Another interesting example is given by the dispute between Euler and d'Alembert around the problem of string vibration, which occurred the second half of the XVIII century. In 1747, d'Alembert had given a formulation to the wave equation, which is the following:

\[
\frac{\partial^2 y}{\partial t^2}
=
c^2 \frac{\partial^2 y}{\partial x^2}
\]

This equation has the solution of the form

\[
y(x,t) = f(x-ct) + g(x+ct),
\]

being \(f\) and \(g\) suitable functions. The problem concerned the class of functions from which \(f\) and \(g\) should be drawn in order to capture all solutions of the equation. D'Alembert had in mind analytic functions represented by smooth formulas, whereas Euler adopted a broader conception of function, one that was not necessarily expressible by a single formula. Prior to this controversy, the notion of a real-valued function was closely tied to its representation by an analytic expression, and Euler himself---see \textcite[p.~117]{Maddy1997}---often characterized functions in these terms.

However, the wave equation led Euler to revise this view. The initial conditions of this equation were physically interpreted as describing the shape of a string \(y(x,0)\) in space, which could, at least in principle, assume a serrated and non-smooth form. Such a shape might be non-differentiable or even fail to be representable by a single analytic formula, a possibility that D'Alembert rejected.\footnote{Such serrated functions have played an important role in the subsequent history of analysis. One of the most paradigmatic examples is the Weierstrass function, formally introduced in 1872. This function is defined by the infinite series \(\sum_{n=0}^{\infty} a^n \cos(b^n \pi x)\), where \(a \in (0,1)\), \(b\) is an odd positive integer, and \(ab > 1 + \frac{3}{2}\pi\).}

Euler maintained that functions ``are equally satisfactory, whether they are expressed by some equation or whether they are traced in any fashion, in such a way as not to be subject to any equation'' (quoted and translated in \textcite[p.~288]{Truesdell1960}). I do not think Euler was thinking on a set $\mathbb{R}^\mathbb{R}$ containing all the arbitrary real-valued functions in $\mathbb{R}$. Euler's contribution was not the introduction of a precise new structure, but the recognition that the concept of function should be broadened beyond analytic expressions. In modern terms, this corresponds to a transition from a restricted class of functions to progressively richer function spaces.

This conceptual transition proved extraordinarily fruitful, as it enabled the mathematical treatment of a much broader range of physical and analytical phenomena. By allowing functions that were not necessarily smooth or representable by a single analytic expression, mathematicians could rigorously formulate and solve problems arising from wave propagation, heat diffusion, and, later, quantum mechanics. Beyond its problem-solving power, this broader conception of function also fostered a remarkable explanatory unification: seemingly disparate areas such as geometry, analysis, and spectral theory came to be understood within a common framework, particularly through the development of Hilbert-space methods. Moreover, the expanded notion of function exhibited considerable generative capacity, giving rise to entire new branches of mathematics, including functional analysis, operator theory, the modern theory of partial differential equations, and harmonic analysis. What began as a dispute over the admissible shapes of vibrating strings thus contributed to a profound reconfiguration of the mathematical landscape. Modern structures such as Hilbert spaces and function spaces like $(L^2(\mathbb R))$ may be viewed as particularly striking outcomes of this process. Their acceptance did not arise from an antecedent adoption of those structures themselves, but from their sustained success in extending and organizing mathematical practice. In the terminology introduced in Section \ref{quasi-emp-structuralism}, such structures were incorporated because of their internal fruitfulness with respect to simpler structures, such as $(\mathbb R,C^2(\mathbb R))$, which are closer to the conception of function defended by d'Alembert.

I want to make clear that not all legitimacy disputes in mathematics are decided in favor of the entity under consideration. A particularly instructive example is provided by the history of infinitesimals. Originally, in the expression

$$\int_a^b f(x),dx$$

the symbol (dx) was not merely regarded as part of the notation but was often understood as denoting an infinitesimal quantity in its own right. Likewise, operations such as

$$d(f \cdot g)=df \cdot g + f \cdot dg$$

were accepted and widely employed throughout the seventeenth and eighteenth centuries.\footnote{See, for instance, \textcite{Bos1974} for a historical account of infinitesimals. The topic is also discussed by \textcite{Lakatos1976}, who emphasizes the important role that conceptually problematic entities may nevertheless play in mathematical practice. }

Later, however, infinitesimals were largely displaced by a more fruitful framework, namely the limit concept that emerged during the nineteenth-century formalization of calculus. This framework offered greater explanatory unity and rigor, while integrating more naturally into the structure of the real line. From the perspective of quasi-empirical structuralism, the point is not that infinitesimals were devoid of mathematical value. On the contrary, they played an important heuristic and developmental role. Rather, the relevant observation is that the limit framework ultimately proved more fruitful according to the criteria introduced in Section 3 and therefore became the dominant foundation for analysis. It is true that expressions involving $dx$ receive rigorous interpretations in later theories, such as differential forms in geometry, though not in precisely the same sense originally intended. The point I would like to make is not that infinitesimals were meaningless, but that the original conception lost its central role.

It is true that contemporary mathematics contains alternative frameworks, such as the surreal numbers and nonstandard analysis, in which infinitesimal-like entities can once again be rigorously defined.\footnote{See \textcite{Knuth1974} for an accessible introduction to surreal numbers. I do not deny the conceptual interest or mathematical depth of such frameworks. The point is simply that they have not achieved the same degree of integration into mainstream mathematical practice as the classical limit-based approach.} Their existence nevertheless reinforces the central lesson of this section: adoption of mathematical structures is not determined solely by logical consistency or conceptual possibility, but also by the extent to which a structure proves fruitful within mathematical practice.

The examples discussed in this section exhibit a common pattern. Mathematical structures are not typically adopted simply because they are logically consistent or formally definable. Rather, they gain legitimacy through their capacity to solve problems, unify previously disconnected domains, and generate new lines of inquiry. Complex numbers, modern function spaces, and other successful mathematical structures earned their place within mathematics because they proved fruitful according to these criteria. Conversely, structures or concepts that fail to play such a role may remain mathematically coherent without becoming central to mathematical practice. From the perspective of quasi-empirical structuralism, the significance of a mathematical structure is therefore inseparable from its position within a broader network of applications, explanations, and developments.

%=========================
\section{The natural numbers and the real line: structures determined by practice?}
\label{natural-reals}

In this section I return to the question of the determinacy of mathematical structures, introduced in Section \ref{structuralims-ontology-epistemology} as a distinction between adoption of structures and adoption of axiom systems. I argue that structures such as $(\mathbb{N}, 0, S)$ and $(\mathbb{R},1, +, \cdot, <)$ have proved fruitful in mathematical practice precisely because mathematicians have treated them as determinate structures rather than merely as open-ended collections of models satisfying an axiom system. This point is reinforced by the distinctive role these structures play beyond pure mathematics. Their deep entrenchment in scientific applications and everyday forms of reasoning gives them a stability and significance not shared by many more specialised mathematical structures. Counting, measurement, spatial representation, and quantitative modelling all rely on them in essential ways. In the terminology introduced in Section \ref{quasi-emp-structuralism}, they therefore provide some of the clearest examples of externally applicable structures.

%---------Natural numbers-----------------
$\mathbb{N}$ is the simplest infinite mathematical structure, yet it is remarkably rich, in many cases enough for most applications in physics and many areas of mathematics. Its applications are so widespread that they are difficult to survey exhaustively. As \textcite[10]{maclane1986mathematics} observes, ``arithmetic operations were invented (or discovered?) because they have all manner of practical uses in financial or scientific calculations.'' Practical activities have also motivated the development of more sophisticated combinatorial concepts grounded in the natural numbers. Mac Lane illustrates this with the example of gambling, where questions about the probabilities of different card arrangements naturally lead to the study of finite permutation groups:

\begin{adjustwidth}{1cm}{1cm}
\begin{quote}
  \small
Such counts are useful in gambling or speculating. Choose three cards in succession from an (ordered) deck of thirteen; what is the chance that they come out in a direct or reverse order? It is the ratio of favorable cases [(123) or (321)] to the total number 6 of cases (of permutations). \parencite[12]{maclane1986mathematics}
\end{quote}
\end{adjustwidth}

The central challenge, however, is to explain how such applications can support the claim that arithmetic is determinate. A standard response appeals to categoricity results for arithmetic. Since most mathematical and scientific applications of arithmetic rely on the axioms of PA, a categorical characterization of those axioms would secure a unique structure as their intended model. In that case, arithmetical truth could be identified with truth in that structure. Dedekind's proof that the second-order version of PA is categorical \parencite{Dedekind1888} is often taken to provide precisely such a characterization. This result has been widely accepted in mathematical practice as capturing the structure of the natural numbers; see, for example, \textcite[47]{maclane1986mathematics} and \textcite[141]{Ferreiros2023ConStructuralism}.

\textcite[19]{Feferman2011a} and \textcite[511]{Koellner2016}, fundamentally because it is proven within $\mathrm{ZFC}$ (or perhaps in weaker versions such as $\mathrm{ZF}^- -P + Pow(\omega)$)\footnote{Here $Pow(\omega)$ stands for the axioms asserting the existence of the powerset of the set of all natural numbers.} , while these background set theories are themselves non-categorical. As a result, the categoricity theorem secures uniqueness only relative to a particular model of the underlying set theory. Although any two realizations of the natural-number structure within a given model are isomorphic, different models of the background theory may contain different realizations of that structure. Consequently, the categoricity result alone does not establish an absolutely unique natural-number structure. 

For this reason, I am reluctant to appeal to categoricity results in full second-order logic as a sufficient basis for the determinacy of arithmetical statements. This echoes the concerns raised in Section \ref{structuralims-ontology-epistemology}, where I expressed reservations about \textcite{shapiro1991foundations}' reliance on full second-order logic in accounting for the determinacy of mathematical structures.

I believe that many of these difficulties can be resolved by taking seriously the distinction between adoption of structures and adoption of axiom systems. In particular, the role that mathematical practice assigns to results such as Dedekind's categoricity theorem can only be adequately understood in light of this distinction. Mathematicians do not typically regard the theorem as establishing the determinacy of arithmetic from a neutral, external standpoint. Rather, it is understood and proved \textit{within} mathematical structures that are already taken as sufficiently determinate for the purposes of the proof. In particular, the categoricity theorem for $\mathrm{PA}_2$ is established within structures incorporating $\mathbb{N}$ and $\mathcal{P}(\mathbb{N}) \cong \mathbb{R}$.

From this perspective, the theorem should be seen not as grounding the determinacy of arithmetic, but as expressing a structural relationship that holds within a broader mathematical framework already committed to those structures. And the reason why mathematicians adopting these structures is, in my view, quite related to practice, in the spirit of quasi-empirical structuralism. The remaining question, then, is why mathematicians are justified in treating those background structures as sufficiently determinate. My suggestion is that the answer lies largely in mathematical practice itself. In the spirit of quasi-empirical structuralism, adoption of structures such as $\mathbb{N}$ and $\mathbb{R}$ is not primarily secured by foundational arguments, but by their sustained success across a wide range of mathematical, scientific, and everyday applications.

This is something that, for instance, non-standard models of first-order arithmetic fail to do. Such a model $(M, 0_M, S_M)$ is characterized by the incorporation of a non-standard element, an element $a\in M$ such that $0_M, S_M(0), S_M(S_M(0)),... < a$. As the property of being a non-standard element is not expresible in first-order logic, then first-order $\mathrm{PA}$ cannot exclude the appearance of these elements. Any enumerable non-standard model of arithmetic has the following order type:

$$\omega+ \mathbb{Z} \eta$$ 

where $\eta$ represents the usual order type of the rational numbers. Thus, after an initial segment isomorphic to the standard natural numbers, the model contains infinitely many copies of the integers arranged densely like the rationals.

These non-standard models models of $PA$ fail to capture the notion of counting and indefinite iteration as the standard structure $(\mathbb{N},0,S)$ does. In particular, many applications presuppose a one-to-one correspondence between the natural numbers and processes generated by repeated succession. A simple example is provided by the sequence of future days in a calendar. Each future day can naturally be associated with a unique natural number representing the number of days elapsed from a given starting point. The standard model of arithmetic captures this correspondence exactly. By contrast, non-standard models contain additional elements that do not correspond to any finite stage of such a process, and therefore cannot be interpreted as representing the structure underlying ordinary counting practices.

%-----------Real numbers-----------------
A similar point applies to non-standard models of theories intended to characterize the real line. Although such models may satisfy the relevant axioms, they do not play the same role in spatial representation, measurement, and scientific modelling as $(\mathbb{R},1, +, \cdot, <)$. This structure is commonly characterized as a complete ordered field satisfying the least-upper-bound principle. Under suitable higher-order formulations, these axioms characterize the real numbers uniquely up to isomorphism, much as second-order Peano Arithmetic characterizes the natural number structure. As before, however, the resulting categoricity depends on assumptions that go beyond first-order logic.

Similarly to arithmetic, restricting attention to first-order theories intended to characterize the real line leads to the appearance of non-standard models. A classical example is provided by the field of hyperreal numbers $^*\mathbb{R}$, which contains all real numbers together with infinitesimal elements $\varepsilon$ such that, for every standard natural number $n$:\footnote{See \textcite{robinson1966nonstandard} for more insights on non-standard models of analysis.}

$$ 0 < \varepsilon < \frac{1}{n} $$

These non-standard structures are nevertheless treated quite differently in mathematical practice from the standard real line. The latter provides the framework within which continuity, magnitude, and geometric space are successfully represented and studied. Its privileged status is closely connected to the role that the real numbers play in a wide range of human and scientific activities. In this spirit, \textcite[93]{maclane1986mathematics} emphasizes the central place of the real numbers within mathematics:

\begin{adjustwidth}{1cm}{1cm}
\small
\begin{quote}
[The system of real] numbers form the most central structure in all of Mathematics. Like other structures it arises from more primitive human activities-in this case, from the measurement and comparison of magnitudes of various kinds.
\end{quote}
\end{adjustwidth}

A natural objection to the position developed here comes from constructivist approaches to mathematics, such as the ones given in \textcite{Bishop1967} or \textcite{bridges1987varieties}. Constructivists typically require mathematical objects to be given through explicit methods of construction and are therefore often reluctant to accept completed infinite totalities whose elements cannot, even in principle, be individually generated. Applied to the real numbers, this leads many constructivist accounts to restrict attention to classes of reals that admit some effective or computational description. Such approaches have considerable philosophical appeal, particularly from the perspective of mathematical practice, since they seek to ground mathematical meaning in concrete procedures rather than abstract existence assumptions. Nevertheless, one may question whether such restrictions fully capture the role that the real numbers play in geometry, measurement, and scientific modelling. This concern is expressed by Mac Lane:

\begin{adjustwidth}{1cm}{1cm}
\small
\begin{quote}
Other approaches restrict attention to those real numbers which are explicitly “constructible” in some specific sense, perhaps by means of recursive functions of natural numbers (...). In my view, such requirements of constructivity tend to emphasize arithmetic at the expense of geometry. Since the real numbers should subsume both arithmetic and geometric aspects of mensuration, these practical uses seem better formalized by the standard axioms for the field of real numbers. \textcite[107]{maclane1986mathematics}
\end{quote}
\end{adjustwidth}

The discussion of arithmetic and the real line suggests that determinacy in mathematics cannot be explained solely by appeal to categoricity results. Although such results are often taken to characterize the intended structures of arithmetic and analysis, they themselves presuppose background frameworks whose determinacy is open to further philosophical scrutiny. What ultimately distinguishes the natural numbers and the real line is not merely that they can be characterized categorically, but that they occupy a uniquely entrenched position within mathematical, scientific, and everyday practice. Once such structures are adopted, questions that remain unresolved at the level of a first-order axiom system—whether because they are undecidable from the axioms or because they depend on features varying across models—receive definite answers, at least in principle, within a privileged structure. From the perspective of quasi-empirical structuralism, this privileged status is closely connected to their role as externally applicable structures, deeply embedded in practices of counting, measurement, spatial representation, and scientific modelling.\footnote{This does not entail that non-standard models can never be fruitfully employed in mathematics. For example, hyperreal fields have been used in non-standard analysis and asymptotic methods; see \textcite{lightstone1975nonarchimedean}. The claim is only that such structures do not occupy the same foundational position as the standard real line in practices of measurement, geometry, and scientific modelling.}

The question, then, is whether a similar argument can be given for the structures associated with contemporary set theory. Do mathematicians adopt the cumulative hierarchy in the same sense that they adopt the natural numbers and the real line, or does set theory occupy a different position within mathematical practice? This is the issue to which I now turn.

%===========================
\section{The approach to set theory}
\label{set-theory}

%-------Intro: set theory is not externally fruitful-----------

Most people would agree that abstract set theoretical axiom systems, such as $\mathrm{ZFC}$, are not externally applicable. \textcite{Feferman1993} argued that really weak axiom systems, even $\mathrm{PA}$ could suffice for the applications of mathematics in physics. For these reason, in case mathamaticians have adopted abstract set-theoretic structures or axiom systems, that is expected to have happened due to their internal fruitfulness, that is, its utility within other areas of mathematics. In this regard, I think there are three fundamental reason to argue the internal fruitfulness of such abstract set theoretic structures:
 
\begin{enumerate}

\item \label{ext-axioms} Axioms in set theory serve as an extension of weaker theories that cam be applied to previously adopted structures, such as $\mathbb{N}$ or $\mathbb{R}$.

\item \label{common-language} Set theory has served as a common language for other areas of mathematics, allowing disparate concepts, reasoning and results to be expresed in a unifying language.

\item \label{metamathematics} Set theory has enabled metamathematical considerations, i.e., has been useful to determine the scope and reach of previously adopted axiom systems.

\end{enumerate}

In this Section I argue that for any of those uses, differently to arithmetic and geometry, mathematicians have focused in set theory in its axiomatic form, with no need of an underlying privileged structure.

%----First use: extension of the axioms----------------
Regarding \ref{ext-axioms}, working in the language of set theory, one may consider formulas whose quantifiers are all relativized to the set $\omega$ of natural numbers. Such formulas will be called \emph{arithmetical sentences}. We then define

$$ \mathrm{Arith}(\mathrm{ZFC}) = \{\varphi : \varphi \text{ is arithmetical and } \mathrm{ZFC}\vdash\varphi\}.$$

Thus, $\mathrm{Arith}(\mathrm{ZFC})$ is the collection of all arithmetical consequences of $\mathrm{ZFC}$. This theory enjoys several basic properties. First, $\operatorname{Th}(\mathrm{PA}) \subseteq \mathrm{Arith}(\mathrm{ZFC})$, since every theorem of first-order Peano Arithmetic is provable in $\mathrm{ZFC}$ under the standard interpretation of arithmetic in set theory. Furthermore,

$$ \mathrm{Arith}(\mathrm{ZFC}) \subseteq \operatorname{Th}(\mathbb{N}), $$

where $\operatorname{Th}(\mathbb{N})$ denotes the complete first-order theory of the standard model of arithmetic. In other words, assuming that $\mathrm{ZFC}$ is consistent, every arithmetical sentence provable in $\mathrm{ZFC}$ is true in $\mathbb{N}$.

Similarly, one may consider the consequences of $\mathrm{ZFC}$ obtained by restricting quantification to subsets of $\omega$, that canonically code real numbers. Assuming that $\mathrm{ZFC}$ is consistent, every such consequence is true in the standard structure of second-order arithmetic $(\omega, \mathcal{P}(\omega))$, and hence belongs to $\operatorname{Th}((\omega, \mathcal{P}(\omega)))$. Since $\mathcal{P}(\omega)$ is in bijection with $\mathbb{R}$, this theory captures the consequences of $\mathrm{ZFC}$ concerning the real line, i.e. the structure $(\mathrm{R}, 1, +, \cdot, <)$.

For this reason, the adoption of standard structures such as $\mathbb{N}$ and $\mathbb{R}$ makes the use of $\mathrm{ZFC}$ particularly interesting. Although formulated as a theory of sets, $\mathrm{ZFC}$ has consequences concerning these lower-order structures that go beyond what can be proved in the usual first-order axiomatizations of arithmetic and analysis. A notable example is Goodstein's theorem, an arithmetical statement that is unprovable in first-order Peano Arithmetic but can be established using methods involving transfinite ordinal induction, which can be formalized within $\mathrm{ZFC}$ (indeed, considerably less strength suffices). This illustrates that it is not unnatural to regard stronger set-theoretic principles as yielding genuine information about more elementary mathematical structures.


%--------Reason 2: unifying language------------
Point \ref{common-language} is closely related to this. Perhaps the greatest success of set theory has been its integration as a common language and conceptual framework for virtually all of mathematics. The notion of set provides a uniform way of representing mathematical objects, functions, spaces, relations, and structures from widely different areas. In this sense, set theory exhibits considerable explanatory fruitfulness by supplying a shared language through which mathematical theories can be formulated and compared.

I would argue, however, that this success is primarily linguistic and conceptual rather than foundational. Their success in mathematical practice therefore provides reasons for adopting set-theoretic modes of reasoning, but not necessarily for privileging a specific cumulative hierarchy as the intended universe of sets.

This point can be illustrated by considering individual axioms. Principles such as the Axiom of Choice are frequently employed in mathematical practice, but it is often unclear whether their full strength is required, as opposed to weaker principles such as Dependent Choice. An even more revealing example is provided by Replacement. In practice, the full Axiom of Replacement is almost always replaced by its weaker bounded\footnote{The Axiom of Replacement asserts the existence of the set ${\varphi(z)\mid z\in x}$, for any set $x$ and any univocal formula $\varphi$ in the language. The bounded version requires all quantifiers occurring in $\varphi$ to be relativized to previously given sets, that is, to be of the form $\forall u\in t$ or $\exists v\in s$.} version---see for instance \textcite[postnote]{maclane1986mathematics} or \textcite[2]{McLarty2020}. The fact that ordinary mathematics can typically be carried out using such weaker principles suggests that the practical utility of set theory is not tightly connected to commitment to the full cumulative hierarchy envisioned by ZFC.

A similar observation applies to the Power Set Axiom. At first sight, it seems built into the very definition of a topological space, since a topology on a set is defined as a collection of subsets of that set. Yet the working topologist typically reasons schematically about a particular family of subsets satisfying certain closure conditions, rather than about a completed totality of all subsets. Theorems are proved about the given topological structure itself, and not about the full power set from which the topology is drawn. Such schematic reasoning differs significantly from commitment to an independently existing power set in the full set-theoretic sense.

It is true, for instance, that the definition of the discrete topology, It is true, for instance, that the definition of the discrete topology, where the topology consists of the entire power set, appears to invoke the full Power Set Axiom. Here, however, it is important to distinguish two different situations. On the one hand, a mathematician may reason directly about a particular topological structure. This is closer to the notion of adopting a mathematical structure introduced in Section \ref{structuralims-ontology-epistemology}. On the other hand, the existence of a discrete topology may be invoked only schematically, for example in the construction of counterexamples or in the formulation of general results about topological spaces. In the former case, the structures actually manipulated in practice are very often countable and explicitly given. In the latter case, one merely appeals to a broader set-theoretic framework without thereby committing oneself to any particular cumulative hierarchy. The mathematical usefulness of such reasoning therefore does not by itself provide grounds for privileging one set-theoretic universe over its competitors.

These examples are necessarily selective and are intended only to illustrate a broader tendency within mathematical practice. Their significance for the present discussion is that they suggest a distinction between adopting set-theoretic methods and adopting a privileged set-theoretic universe. The central claim of this paper is that the former is widespread in contemporary mathematics, whereas the latter is considerably less evident.

%---Reason 3: metamathematics-------------
And finally, regarding point \ref{metamathematics}, axiomatic set theory often functions in a metamathematical role, providing results about the proving power of specific axiom systems for formalized mathematical languages. Indeed, the connection between mathematical provability and formal derivability from the axioms of set theory is so close in contemporary practice that \textcite[373]{maclane1986mathematics} proposes the following standard of rigor:

\begin{adjustwidth}{1cm}{1cm}
\small
\begin{quote}
Proof in Mathematics is both a means to understand why some result holds and a way to achieve precision. As to precision, we have now stated an absolute standard of rigor: A Mathematical proof is rigorous when it is (or could be) written out in the first order predicate language $\mathcal{L}(\in)$ as a sequence of inferences from the axioms ZFC, each inference made according to one of the stated rules.
\end{quote}
\end{adjustwidth}

A paradigmatic example of metamathematical reasoning about set theory is provided by independence results. These results establish that certain statements—most notably the Continuum Hypothesis—cannot be proved or refuted from the axioms of a theory such as ZFC. Crucially, such arguments do not presuppose a uniquely determined set-theoretic universe of the kind whose status is under investigation. Rather, they proceed by reasoning about formal theories, proofs, and models from an external framework known as the metatheory. In many presentations, the metatheory is considerably weaker than the set-theoretic theory being studied and requires only comparatively modest assumptions concerning formulas, derivations, and formal reasoning. As \textcite[11]{Kunen2011} observes: 

\begin{adjustwidth}{1cm}{1cm}
\begin{quote}
  \small
All we assert as being ‘really true’ is elementary finitistic reasoning: this is called the metatheory. Using this reasoning, we can explain what a logical formula is and what a formal proof is (...). Then, all discussions about infinite and uncountable sets take place within ZFC, which is called the formal theory.
\end{quote}
\end{adjustwidth}

Returning to the point under discussion, it is interesting to note that these applications make no appeal to a standard or privileged model of set theory. Indeed, many metamathematical applications focus primarily on formal logic, and the models employed are often deliberately chosen to be non-standard:

\begin{adjustwidth}{1cm}{1cm}
\begin{quote}
\small
Take a case like Cohen's method of forcing where you build a non-standard model (...). This is done within set theory and it is revealed in the first step that one is dealing with a non-standard model since it is a countable object. In the second approach one builds class size models (...) but one builds a Boolean-valued model and, once again, it is immediately revealed that it is non-standard. \parencite[p. 16]{Koellner2022}
\end{quote}
\end{adjustwidth}

This feature becomes particularly apparent in forcing and independence proofs. As Koellner observes, the models used in such arguments are typically recognized from the outset as non-standard, whether because they are countable transitive models or because they arise through Boolean-valued constructions. This observation is not merely technical; independence proofs proceed by reasoning about formal theories and models without requiring access to a uniquely determined universe of sets. While set theorists possess a variety of criteria for recognizing models as non-standard, no comparably straightforward criterion exists for identifying a standard model. Consequently, the success of forcing arguments demonstrates the fruitfulness of set-theoretic methods while remaining neutral on the existence of a privileged set-theoretic structure.

%-------Conclusion: adoption of set theory but not in structural form-------------------------
%-------Conclusion: adoption of ZFC but not stronger axioms that could decide CH--------------
It is interesting to note that, unlike many other cases in mathematics, the development of the axioms of set theory preceded the articulation of the structure modelling them, the cumulative hierarchy. The hierarchical conception of the set-theoretic universe emerged later in the work of \textcite{vonNeumann1925} and \textcite{Zermelo1930}, largely in order to explain how the axioms could be understood, to avoid the classical paradoxes, and to provide a structured conception of the universe of sets. In this respect, the situation slighly differs from cases such as arithmetic, where the underlying structure was already integrated in mathematical practice before its formal axiomatization. One might therefore say that, in set theory, the structure was introduced in order to make sense of the axioms, rather than the axioms being introduced to describe a previously adopted structure. Indeed, understanding the cumulative hierarchy independently of the formal axioms is not straightforward:

\begin{adjustwidth}{1cm}{1cm}
\small
\begin{quote}
Within ZFC set theory, the cumulative hierarchy can be formally constructed. Before axioms are at hand, the hierarchy is a Platonic myth, clearly visible only to those with a sixth sense for sets. \textcite[373]{maclane1986mathematics}
\end{quote}
\end{adjustwidth}

To sum up, the axioms of set theory prove fruitful with respect to previously adopted mathematical structures, both by extending their theory and by providing a common framework for mathematical reasoning. Thus, regarding the examples discussed in this section, we may reasonably say that

$$(\mathbb N,0,S), \; (\mathbb{R},1, +, \cdot, <) \; \leadsto \; \mathrm{ZFC}.$$

What does not follow, however, is a corresponding adoption of a determinate cumulative hierarchy. In particular, the preceding considerations do not establish

$$(\mathbb N,0,S), \; (\mathbb{R},1, +, \cdot, <) \; \leadsto \; \big( \bigcup_{\alpha}V_\alpha,\in \big).$$

The examples examined above provide reasons for adopting the axioms of set theory as fruitful mathematical tools, but not necessarily for adopting a privileged set-theoretic universe as a determinate mathematical structure.
 
%-----------References to set-theorists-------
I do not think this perspective undermines the fruitfulness of set theory, nor should it be rejected by defenders of set-theoretic practice. For instance, \textcite{maddy1988believing1,maddy1988believing2} argues that set-theoretic axioms acquire justification primarily through extrinsic rather than intrinsic considerations. Her account provides a compelling explanation of how increasingly strong set-theoretic principles gain legitimacy within mathematical practice. Moreover, the criteria she invokes are largely compatible with the quasi-empirical considerations discussed above. In particular, Maddy's notion of extrinsic justification concerns the adoption of axioms rather than the adoption of structures. What remains less clear, however, is whether the fruitfulness of an axiom system is sufficient to establish the existence of a uniquely determined set-theoretic structure corresponding to it. Maddy's considerations provide reasons to accept and employ set-theoretic axioms, but the further transition from axiomatic fruitfulness to commitment to a privileged cumulative hierarchy requires additional argument.

\section{How \texorpdfstring{$\mathrm{CH}$}{CH} fits in quasi-empirical structuralism}
\label{section-CH-case}

We may now return to the question that motivated this paper: whether the Continuum Hypothesis should be regarded as a definite mathematical problem. From the perspective of quasi-empirical structuralism, there are two possible routes by which such definiteness might be secured. First, mathematical practice might adopt an axiom system that settles $\mathrm{CH}$ and becomes sufficiently entrenched to form part of the background framework of mathematical reasoning. Second, mathematical practice might adopt a determinate structure in whose language $\mathrm{CH}$ is evaluated, thereby fixing its truth value independently of any particular axiomatic extension.

%-----Further axioms----
Regarding the first of these possibilities, we saw in the previous section that mathematical practice consistently relies on axiom systems such as $\mathrm{ZFC}$, or perhaps weaker fragments, both as a common language for mathematics and as a useful metamathematical framework. In the sense introduced above, such systems may therefore be regarded as adopted by mathematical practice. This fact alone, however, is insufficient to secure the definiteness of $\mathrm{CH}$, since $\mathrm{CH}$ is independent of $\mathrm{ZFC}$.

The existence of axioms such as $\mathrm{V=L}$ and $\mathrm{MM}$ does not undermine the present view.\footnote{Martin's Maximum (MM) is a strong forcing axiom introduced by \textcite{ForemanMagidorShelah1988}. It has substantial combinatorial and structural consequences and decides many questions independent of $\mathrm{ZFC}$, among them $\neg\mathrm{CH}$.} On the contrary, the fact that contemporary set theory continues to treat the universes determined by such principles as competing mathematical possibilities is precisely evidence that mathematical practice has not converged on the adoption of a privileged set-theoretic universe. The fruitfulness of $\mathrm{MM}$ and $\mathrm{V=L}$ is genuine, and it fully explains why they have become important objects of mathematical investigation. What it does not explain is why either principle should be incorporated into the shared background assumptions on which substantial areas of mathematics routinely depend. In the terminology of condition \ref{condition4}, neither principle appears to have achieved the stability of reliance characteristic of genuinely adopted frameworks. So, in our notation, we do not have the following sequence:

$$ (\mathbb{N}, 0, S), \; (\mathbb{R},1, +, \cdot, <) \; \leadsto \; \mathrm{ZFC}+\mathrm{MM} $$

nor 

$$ (\mathbb{N}, 0, S), \; (\mathbb{R},1, +, \cdot, <) \; \leadsto \; \mathrm{ZFC}+\mathrm{(V=L)}. $$

%----------Finite-order arithmetic: Weil conjectures, Fermat's last theorem and MC systems----------
Then it remains to address whether some determined structures have been adopted, in which statements as $\mathrm{CH}$ receive definite truth conditions. Having argued in Section \ref{set-theory} that this is not the case for the set-theoretic universe, we must now examine whether a positive answer might nevertheless emerge from weaker frameworks, such as finite-order arithmetic.
A paradigmatic case of the application of stronger foundational frameworks to previously adopted mathematical structures is provided by the \textit{Weil conjectures}. Introduced by Weil in 1949, these consist of four deep conjectures concerning algebraic varieties over finite fields and played a central role in the development of modern algebraic geometry.

What makes the example particularly relevant here is the contrast between the logical complexity of the statements themselves and that of the tools used in their proof. The Weil conjectures can be formulated in comparatively weak systems, such as second-order arithmetic or suitable extensions of Peano Arithmetic. Their eventual proof was obtained by \textcite{Deligne1974}, building on the étale cohomological machinery developed by Grothendieck and his collaborators. In its original formulation, this machinery made use of Grothendieck universes, whose existence requires $\mathrm{ZFC}$ together with a strongly inaccessible cardinal. \textcite{McLarty2020}, however, has shown that the essential mathematical content of these methods can be reconstructed in considerably weaker frameworks, namely finite-order arithmetical systems. These systems, denoted by $\mathrm{MC}$, are, in effect, the weakest within which Grothendieck's methods can be formalized and Deligne's proof carried out.

The systems $\mathrm{MC}$ require the existence of the iterated power sets $P^n(\mathbb N)$ for each natural number $n$. In this respect, they resemble the higher-order frameworks discussed in Section \ref{set-theory}. Yet the existence of such iterated power sets does not by itself determine a privileged cumulative hierarchy. At most, it licenses the use of specific higher-order structures required for particular mathematical purposes. Consequently, the considerations advanced in Section \ref{set-theory} apply equally here: the successful use of these frameworks provides strong reasons to employ them, but considerably weaker reasons to regard the corresponding higher-order structures as determinate mathematical objects.

%-------------Fermat's Last Theorem--------------
A similar phenomenon can be observed in the proof of Fermat's Last Theorem. Although the theorem is a statement about the natural numbers, its proof relies on sophisticated machinery from arithmetic geometry and the theory of modular forms. Current foundational reconstructions place this machinery within finite-order arithmetic, though it remains an open question how weak the underlying framework can ultimately be made \parencite[17]{McLarty2020}. 

\textcite{Feferman1999} defended the view that Fermat's Last Theorem should ultimately be provable within substantially weaker systems, perhaps even within $\mathrm{PA}$. Whether or not this expectation is correct is largely irrelevant to the present discussion. My argument does not require that ordinary mathematics reduce to first-order arithmetic, only that the structures actually relied upon by mathematical practice fall short of determining a privileged universe of sets.


\section{Conclusion}
\label{Conclusion-chapter2}

In any case, if there is a boundary between adopted mathematical structures and merely available foundational frameworks, it appears to lie somewhere within finite-order arithmetic rather than at the level of the full cumulative hierarchy. This conclusion should not be understood as denying the mathematical legitimacy of stronger frameworks. Contemporary set theory has produced a rich variety of fruitful structures whose investigation has substantially deepened our understanding of mathematics. The present claim is only that such structures have not been adopted in the same sense as the natural numbers, the real line, or the finite-order frameworks on which ordinary mathematical practice routinely relies.

Whether this reflects an inherent limitation of our mathematical capacities or merely the current state of mathematical development remains an open question. Nevertheless, a common theme running through the perspectives considered in this paper—from Wittgenstein's emphasis on use, through Lakatos's and Mac Lane's focus on mathematical practice, to Feferman's conceptual structuralism—is that mathematical meaning and determinacy arise from structures that human beings can successfully employ, understand, and integrate into their mathematical activity. In this regard, I share, at least in spirit, although not necessarily in its specific content, the following words of \textcite[p. 2]{Bishop1967}:

\begin{adjustwidth}{1cm}{1cm}
\small
\begin{quote}
Mathematics belongs to man, not to God. We are not interested in properties of the positive integers that have no descriptive meaning for finite man (\dots). If God has mathematics of its own that has to be done, let him do it himself. 
\end{quote}
\end{adjustwidth}

Of course, judgments of mathematical fruitfulness are neither fixed nor timeless. Mathematical practice evolves, and structures that presently occupy a peripheral role may in the future acquire a central position within the discipline. Nothing in the present account excludes this possibility. On the contrary, the quasi-empirical perspective developed here is intended to be dynamic rather than static. The adoption of mathematical structures is understood as an ongoing historical process, one that depends on their continued integration into mathematical practice and their capacity to support further developments. Whether future mathematical developments will lead to the adoption of stronger set-theoretic structures remains an open question.

The claim defended in this paper is only that contemporary mathematical practice has not yet done so. More importantly, the issue cannot be settled merely by appealing to proof-theoretic considerations. The point is not merely that existing proof methods are insufficient to decide CH. Even the background structures and axiom systems currently adopted in mathematical practice fail to determine a unique answer. Consequently, there is at present insufficient reason to regard the Continuum Hypothesis as possessing a uniquely determined truth value.


%-----------Tercer capítulo-------------
\chapter{Mathematics as institutions, and how this affect the case around \texorpdfstring{$\mathrm{CH}$}{CH}}
\label{chapter3}

%-----------Conclusión-------------
\chapter{Conclusion: mathematics, conventions and the Continuum Hypothesis}

Perhaps externally applicable structures are more likely to be determined. Although in an idealized way, these structures are related to proccesses or phenomena in real-life experience that our intuition considers determinate. 

Mathematicians adopt the natural numbers+ there is no recursively enumerable set of axioms that determines this set= there can be cases of adopting a structure that cannot be reduced to a set of axioms.

If $L$ is a large cardinal axiom, the theories $ZFC+L$ and $PA + \bigcup_{n < \omega} Con((ZFC + L)\upharpoonright n)$ are mutually interpretable, being the latter one framed in the language of arithmetic. This means that the logical consequences of $ZFC + L$ in the language of set theory can be understood as arithmetical sentences through a well chosen bijection. That fits well with the formalist approach to the metamathematical results, where set theoretical sentences and deductions are understood as the object of study, while the core of the reasoning is done in the metatheory of formal finitary logic. This in fact means that, even if accepting the fruitfulness of set theoretical considerations, specially regarding consistency proofs, that reasoning can be reduced to the language of arithmetic. So, in the end, the only mathematical structure one has to adopt in this case is the natural number structure, without need of a freestanding and platonic universe of sets.

\section{Some more examples on fruifulness and quasi-empirical structuralism}
This is an example taken from \textcite[139]{Kechris1995} Let $T$ be a nonempty pruned tree on $A$ and let $X \subseteq [T]$ be closed. Thus $X=[S]$ for $S$ a subtree of $T$. Define by transfinite recursion $S_\xi \subseteq T$ as follows: \[ p \in S_0 \iff p=(a_0,\ldots,a_{2n-1}) \in T \setminus S, \] \[ p \in S_{\xi+1} \iff p=(a_0,\ldots,a_{2n-1}) \in T \;\&\; \forall a_{2n} \bigl( (p, a_{2n}) \in T \Rightarrow \exists a_{2n+1} \bigl( (p, a_{2n}, a_{2n+1}) \in S_\xi \bigr) \bigr), \] \[ p \in S_\lambda \iff \exists \xi < \lambda \, (p \in S_\xi), \qquad \text{if $\lambda$ is limit.} \] Do you think this example is interesting? It uses the notion of ordinals in the context of Borel Determinacy?

%---¿Algo de teoría de categorías?----
\section{Perhaps category theory is a more suitable foundation for mathematics}

The axiomatic basis of this theory is the definition of a category $\mathcal{A}$, which consists of the following elements:

\begin{enumerate}
  \item A collection $Obj(\mathcal{A})$, whose members are called the objects of $\mathcal{A}$.
  \item For every pair of objects $A, B \in Obj(\mathcal{A})$, a collection of arrows (or morphisms) from $A$ to $B$. If $f$ is such an arrow, we write $f: A \rightarrow B$.
  \item A composition operation: given two arrows $f: A \rightarrow B$ and $g: B \rightarrow C$, there is an arrow $g \circ f: A \rightarrow C$.
  \item For every object $A$, an identity arrow $id_A: A \rightarrow A$.
\end{enumerate}

These data are required to satisfy the following axioms:

\begin{enumerate}
  \item \textbf{Associativity:} For arrows $f: A \rightarrow B$, $g: B \rightarrow C$, and $h: C \rightarrow D$, we have
  \[
  h \circ (g \circ f) = (h \circ g) \circ f.
  \]
  \item \textbf{Identity:} For every arrow $f: A \rightarrow B$, we have
  \[
  id_B \circ f = f \quad \text{and} \quad f \circ id_A = f.
  \]
\end{enumerate}

\begin{adjustwidth}{1cm}{1cm}
  \small
  \begin{quote}
For these reasons "set" turns out to have many meanings, so that the purported foundation of all of Mathematics upon set theory totters. But there are other possibilities. For example, the membership relation for sets can often be replaced by the composition operation for functions. This leads to an alternative foundation for Mathematics upon categoriesspecifically, on the category of all functions. Now much of Mathematics is dynamic, in that it deals with morphisms of an object L into another object of the same kind. Such morphisms (like functions) form categories, and so the approach via categories fits well with the objective of organizing and understanding Mathematics. That, in truth, should be the goal of a proper philosophy of Mathematics. \textcite[359]{maclane1986mathematics}
  \end{quote}
\end{adjustwidth}

\begin{adjustwidth}{1cm}{1cm}
  \small
  \begin{quote}
Atiyah described mathematics as the “science of analogy.” In this vein, the purview of category theory is mathematical analogy. Category theory provides a cross-disciplinary language for mathematics designed to delineate general phenomena, which enables the transfer of ideas from one area of study to another. The category-theoretic perspective can function as a simplifying1 abstraction, isolating propositions that hold for formal reasons from those whose proofs require techniques particular to a given mathematical discipline. \parencite[ix]{Riehl2014Category}
  \end{quote}
\end{adjustwidth}

Being $(G, +, \cdot)$ any group,  

Possible roles:

evidence for Resnik-style structuralism,
evidence against unique intended structures,
evidence that mathematics often adopts patterns rather than structures,
an alternative foundation.


\section{How does this differenciate from naïve conventionalism?}

\section{Addressing the robustness of the present case around CH}
On the ``dream mathematical solution''.

\section{On set-theoretical pluralism}
On repeating names in mathematics. 
The determinacy of different models in geometry. Applicability of non-euclidean structures.

The problem of translating this parallelism to mathematics. How different set-theoretic notions \textit{must} repeat names.

Indefiniteness does not equiparate to pluralism.

Hamkins argument fundamentally lies upon the parallelism to non-euclidean geometries.

\section{End note}
Why not suppose that the set theoretical universe exists?

\chapter*{More notes and quotes}
\begin{adjustwidth}{1cm}{1cm}
  \small
  \begin{quote}
These considerations may lead us to say that $2^{\aleph_0} > \aleph_0$. That is to say: we can \emph{make} the considerations lead us to that. Or: we can say \emph{this} and give \emph{this} as our reason. But if we do say it---what are we to do next? In what practice is this proposition anchored? It is for the time being a piece of mathematical architecture which hangs in the air, and looks as if it were, let us say, an architrave, but not supported by anything and supporting nothing. \parencite[III, §35]{Wittgenstein1978}
  \end{quote}
\end{adjustwidth}


\begin{adjustwidth}{1cm}{1cm}
  \small
  \begin{quote}
It would be countersensical to exclude his subjective perspective from the investigation of the role that evidence plays in the selection of axioms. Admittedly, this contrasts with the concerns of ‘ordinary science’ where some standard of objectivity is presupposed and subjectivity is regarded as a disturbing element. \textcite[p. 104]{hauser2004}
  \end{quote}
\end{adjustwidth}

Husserl's signitive vs intuitive distintion is crutial in mathematics. The comprehension of these two kind of consciousness acts is what he calls \textit{fulfillment}. Relation to sense vs reference.

\begin{adjustwidth}{1cm}{1cm}
  \small
  \begin{quote}
We experience how in intuition the same objectuality is intuitively realized which was “merely meant” in the symbolic act, and that it becomes intuitable exactly as something determined in such a way as it was initially merely thought (merely meant). \parencite[Part 2, §8]{Husserl1913}
  \end{quote}
\end{adjustwidth}

\printbibliography

\end{document}